\documentclass[11pt]{article}
% Proof development, deliverable A. Phase 1 (skeleton) reviewed 2026-07-12; this file now
% carries phase-2 closures. Thesis-chapter ready: demote \section to \subsection when
% \input into ../thesis/; no package beyond ams* + geometry + hyperref.
% Compile: tectonic skeleton.tex
\usepackage[margin=1.1in]{geometry}
\usepackage{amsmath,amssymb,amsthm}
\usepackage[hidelinks]{hyperref}

\newtheorem{theorem}{Theorem}[section]
\newtheorem{lemma}[theorem]{Lemma}
\newtheorem{corollary}[theorem]{Corollary}
\newtheorem{proposition}[theorem]{Proposition}
\newtheorem{claim}[theorem]{Claim}
\theoremstyle{definition}
\newtheorem{definition}[theorem]{Definition}
\newtheorem{convention}[theorem]{Convention}
\newtheorem{remark}[theorem]{Remark}

% status tags: every lemma carries exactly one, per the anti-fabrication discipline
\newcommand{\CLASSICAL}[1]{\textnormal{\textsc{[classical: #1]}}}
\newcommand{\OWED}{\textnormal{\textsc{[owed]}}}
\newcommand{\MACHINE}{\textnormal{\textsc{[machine]}}}
\newcommand{\OPEN}{\textnormal{\textsc{[open]}}}
\newcommand{\CLOSED}{\textnormal{\textsc{[closed]}}}

\newcommand{\Flags}{\mathrm{Fl}}
\newcommand{\dev}{\mathrm{dev}}
\newcommand{\Aut}{\mathrm{Aut}}
\newcommand{\Stab}{\mathrm{Stab}}
\newcommand{\Sym}{\mathrm{Sym}}
\newcommand{\A}{\mathcal{A}}
\newcommand{\T}{\mathcal{T}}
\newcommand{\R}{\mathbb{R}}
\newcommand{\C}{\mathbb{C}}
\newcommand{\Z}{\mathbb{Z}}
\newcommand{\im}{\mathrm{im}}
\newcommand{\Isom}{\mathrm{Isom}}

\title{Completeness and correctness of \v{C}trn\'act's combinatorial enumeration
of $k$-uniform tilings by regular polygons\\
\large Proof development (deliverable A; phase 2, round 1) --- TilingAtlas working note}
\author{}
\date{July 12, 2026}

\begin{document}
\maketitle

\paragraph{Status (updated after round 3: certificates + audit).}
The phase-1 skeleton was reviewed and approved on 2026-07-12; rounds 1--3 closed on the
same date. \textbf{Every lemma of deliverable A is \textsc{closed}}, and round 3
discharged the machine and audit obligations: A3--A5 independently audited and re-run,
A6 implemented and passing (44 letters pairwise non-isomorphic), P3 passing in both
directions for $k \le 6$, the deliverable-B implementation audit written
(\texttt{audit-deliverable-B.md}: all eight hooks verified sound, seven fix-obligations
recorded, none correctness-critical on valid input), and the historical $k=8/9$ anomaly
root-caused and closed by falsification (audit \S4; hook 9 below). The classical inputs are down to two, both pinned: the
Gr\"unbaum--Shephard species table (Table 2.1.1, pp.~59--61, verified against the copy
in \texttt{resources/papers}) and the Killing--Hopf classification of closed flat
oriented surfaces. \textbf{0 lemmas are \textsc{open}.} Two structural changes came out
of the closure work, both improvements: (1) round 1 corrected the flag dictionary
($\gamma = \sigma_2\sigma_0$, not $\sigma_0$; Remark~\ref{rem:gammafix}); (2) round 2
replaced the planned orbifold-developability route through the C-block by the
\emph{direction bundle} (Remark~\ref{rem:bundle-arch}), the translation-level cover that
\texttt{eu\_develop} literally iterates; all site symmetries unfold combinatorially, the
developed object is a flat torus, and orbifold theory disappears from the proof
entirely. Residual formalization debt is confined to three convention verifications
flagged in Remark~\ref{rem:referee}; each is finite bookkeeping, not a gap in the
argument. Remaining risk now lives in deliverable B (does the C++ implement the abstract
pipeline?) and in the unaudited machine certificates, exactly where a completeness claim
about a program should keep its risk.

\tableofcontents

%======================================================================
\section{Scope, conventions, and the theorem}\label{sec:scope}

\subsection{Objects}

\begin{definition}[tiling, symmetry, uniformity]\label{def:tiling}
A \emph{tiling} in this note is an edge-to-edge tiling $T$ of the Euclidean plane $\R^2$
by regular polygons of edge length $1$: a countable family of closed regular polygonal
tiles with pairwise disjoint interiors whose union is $\R^2$, such that the intersection
of any two tiles is empty, a common vertex, or a common full edge. $V(T)$, $E(T)$, $F(T)$
denote vertices, edges, faces. The \emph{symmetry group} $G(T)$ is the group of all
isometries of $\R^2$ (orientation-reversing ones included) mapping $T$ to itself. $T$ is
\emph{$k$-uniform} iff $G(T)$ has exactly $k$ orbits on $V(T)$.
\end{definition}

\begin{convention}[chirality merge --- definitional, not derived]\label{conv:chiral}
Two tilings are \emph{equivalent} iff some isometry of $\R^2$, orientation-reversing
allowed, maps one to the other. Enumeration is up to this equivalence: a chiral tiling and
its mirror image count once. This matches the convention of OEIS A068599 and the
Galebach/Soto-S\'anchez oracles, and is a settled repo decision. It is an axiom of the
count, not a theorem.
\end{convention}

\begin{convention}[scope axioms]\label{conv:scope}
(i) Unit edge length is a normalization, not a restriction (a tiling by regular polygons of
common edge length rescales). (ii) The polygon palette of the engine is
$\mathcal{P}=\{3,4,6,12\}$; the octagon is excluded from the palette and handled
analytically by Lemmas O1--O4. (iii) $t1002$ denotes the truncated square tiling
$(4.8.8)$, the unique octagon-bearing tiling (Lemma O3); it is re-added to the $k=1$
catalog by hand.
\end{convention}

\begin{definition}[the two tiling classes]\label{def:classes}
For $k\ge 1$ let
\[
\T_k = \{\text{$k$-uniform tilings containing no octagon}\}/\!\sim, \quad
\T_k^{\mathrm{all}} = \{\text{all $k$-uniform tilings}\}/\!\sim,
\]
where $\sim$ is the equivalence of Convention~\ref{conv:chiral}. By Lemma O4, every member
of $\T_k$ uses tiles from $\mathcal{P}$ only and has, after a rotation, all edge
directions in the $30^\circ$ grid; this is how the ``12-direction'' scope enters, as a
consequence of octagon-freeness rather than as an extra hypothesis.
\end{definition}

\begin{remark}[statement hygiene, correcting the brief]\label{rem:hygiene}
The commissioning brief states the $k=1$ case as $P_1 \cong \T_1 \setminus \{t1002\}$
with $\T_1$ already carrying the grid condition. That is redundant: with the grid (or
octagon-free) condition in the definition, $t1002 \notin \T_1$ and the clean statement is
$P_k \cong \T_k$ for \emph{all} $k \ge 1$, with the octagon accounted for separately by
Theorem~\ref{thm:octagon}. This note uses the clean split.
\end{remark}

\subsection{The abstract pipeline}

The theorem is about the abstract algorithm specified in
\texttt{reference/algorithm.txt}; deliverable B (the implementation-faithfulness audit)
separately certifies that the C++ sources (\texttt{eu\_solver}, \texttt{eu\_pruner},
\texttt{eu\_develop}) implement it. The abstract pipeline, over the alphabet $\A$ of
Section~\ref{sec:alphabet}, with target $k \le \texttt{maxnum}$:

\begin{itemize}
\item $\mathrm{Solve}$: the DFS of Section~\ref{sec:search} over partial stub systems;
  emits every closed configuration that passes the rigidity filter
  (Section~\ref{sec:rigidity}) with $1 \le \#\text{vertices} \le \texttt{maxnum}$.
\item $\mathrm{Prune}$: retains exactly one representative per isomorphism class of
  emitted configurations (Section~\ref{sec:rigidity}).
\item $\mathrm{Develop}$: the geometric development map $M \mapsto \dev(M)$
  (Section~\ref{sec:realizability}).
\end{itemize}

Let $P_k$ be the set of pruned outputs with exactly $k$ vertices ($k$ counting vertex
types; for the regular palette every type is counting).

\subsection{The theorem}

\begin{theorem}[main theorem: the engine]\label{thm:main}
Fix $k \ge 1$ and run the pipeline with $\texttt{maxnum} \ge k$. Then:
\begin{itemize}
\item[(T)] \textbf{Termination.} The run halts.
\item[(S)] \textbf{Soundness.} For every $M \in P_k$, $\dev(M)$ is defined and is a
  tiling whose class lies in $\T_k$; in particular its uniformity is exactly $k$.
\item[(C)] \textbf{Completeness.} For every class $[T] \in \T_k$ there is an $M \in P_k$
  with $\dev(M) \sim T$.
\item[(U)] \textbf{Non-redundancy.} The map $M \mapsto [\dev(M)]$ is injective on $P_k$.
\end{itemize}
Equivalently: $\dev$ induces a bijection $P_k \to \T_k$.
\end{theorem}

\begin{theorem}[octagon scope]\label{thm:octagon}
$\T_k^{\mathrm{all}} = \T_k$ for every $k \ge 2$, and
$\T_1^{\mathrm{all}} = \T_1 \sqcup \{[t1002]\}$.
\end{theorem}

\begin{corollary}[catalog completeness]\label{cor:catalog}
For $k \ge 2$ the developed catalog at $k$ is in bijection with all $k$-uniform
edge-to-edge tilings by regular polygons of common edge length, up to isometry with
reflections; for $k = 1$ the catalog together with $t1002$ is.
\end{corollary}

\begin{remark}[evidence is not proof]
The pipeline reproduces A068599 exactly for $k \le 16$, with per-tiling agreement against
the Soto-S\'anchez data where available. Per the anti-fabrication discipline this is
treated as evidence about the implementation and is used nowhere in the proof of
Theorems~\ref{thm:main} and~\ref{thm:octagon}. The one place the oracle legitimately
enters is deliverable B, as a regression gate.
\end{remark}

\subsection{Obligation map}

\begin{center}
\small
\begin{tabular}{lll}
Obligation & Lemmas & Feeds \\
\hline
1 alphabet completeness & A1--A6 & C (no missing letter), U (labels well defined) \\
2 local-rule soundness/necessity & L1, L2, S4 & C (no wrongful prune), S (via C-block) \\
3 search exhaustiveness & T1, S1--S4 & T, C \\
4 rigidity + dedup exactness & R1--R3, P1--P3 & U, C, S \\
5 realizability & C1--C4 & S \\
6 planar--quotient bridge & B0--B3 & C, S \\
scope (octagon, grid) & O1--O4 & Theorem~\ref{thm:octagon} \\
model well-definedness & D1a, D1b & everything \\
\end{tabular}
\end{center}

%======================================================================
\section{The octagon scope lemmas}\label{sec:octagon}

\begin{lemma}[O1: octagon-bearing species]\label{lem:O1} \CLOSED\;
In an edge-to-edge tiling by regular polygons, the multisets of interior angles at a
vertex incident to an octagon, feasible by the angle equation alone, are exactly
$\{4,8,8\}$ and $\{3,8,24\}$.
\end{lemma}

\begin{proof}
Interior angles are $\theta_n = 180 - 360/n$ degrees, $\theta_n \ge 60$, and the angles at
a vertex sum to $360$. With $a \ge 1$ octagons ($\theta_8 = 135$): if $a = 2$ the
remainder is $90$, indivisible into two or more angles $\ge 60$, so it is a single square
angle: $\{4,8,8\}$. If $a \ge 3$ the sum exceeds $360$. If $a = 1$ the remainder is $225$.
As one angle: $225 > 180$, impossible. As two angles $\theta_m + \theta_n = 225$ with
$\theta_m \le \theta_n$ (so $\theta_m \le 112.5$): the regular interior angles $\le 112.5$
are exactly $\theta_3 = 60$, $\theta_4 = 90$, $\theta_5 = 108$ (there is none strictly
between $60$ and $90$, nor strictly between $90$ and $108$), and their complements to
$225$ are $165$, $135$, $117$; of these only $165 = \theta_{24}$ is a regular angle
($135 = \theta_8$ re-enters the $a \ge 2$ case; $117 = 180 - 360/n$ has no integer
solution), giving $\{3,8,24\}$. As three or more angles, each $\ge 60$: the sums
attainable near $225$ from $\{60, 90, 108, 120, \dots\}$ are
$60{+}60{+}60 = 180$, $60{+}60{+}90 = 210$, $60{+}60{+}108 = 228$, $60{+}90{+}90 = 240$,
and four angles give $\ge 240$, so $225$ is never hit. Hence $\{3,8,24\}$ is the only
$a=1$ solution.
\end{proof}

\begin{lemma}[O2: $3.8.24$ does not occur]\label{lem:O2}
\CLOSED\ (anchor: \CLASSICAL{Gr\"unbaum--Shephard \S2.1})\;
No edge-to-edge tiling by regular polygons contains a vertex whose species is $3.8.24$.
\end{lemma}

\begin{proof}
Suppose vertex $v$ carries tiles $\{T_3, T_8, T_{24}\}$. Consider the triangle $T_3$, with
vertices $v = v_0, v_1, v_2$. At $v_0$, the two edges of $T_3$ separate it from $T_8$ and
$T_{24}$ in some order; say edge $v_0v_1$ borders the octagon. At $v_1$ the incident
angles include $T_3$'s $60$ and the octagon's $135$; the remaining $165$ must decompose
into angles $\theta_n$: as a single angle only $\theta_{24} = 165$ fits; as two or more,
$165 - 60 = 105$ is not an angle and three angles exceed $165$. So $v_1$ also has species
$3.8.24$, and the tile across the \emph{other} edge of $T_3$ at $v_1$, namely $v_1v_2$, is
a 24-gon. Repeating at $v_2$: edge $v_2v_0$ borders an octagon. Now each vertex of $T_3$
must see one octagon and one 24-gon across the two incident triangle edges, i.e.\ the
3-cycle of $T_3$'s edges is properly 2-colored (octagon/24-gon) with colors alternating.
A 3-cycle admits no such alternation, and specifically at $v_0$ both incident triangle
edges now border octagons ($v_0v_1$ by choice, $v_2v_0$ by propagation), contradicting
the species $3.8.24$ at $v_0$. (G\&S state the exclusion of all six double-starred
species, $3.8.24$ included, in \S2.1, p.~59 with Table 2.1.1; pinned against the copy in
\texttt{resources/papers}. The proof above is self-contained regardless.)
\end{proof}

\begin{lemma}[O3: an octagon forces $t1002$]\label{lem:O3} \CLOSED\;
Any edge-to-edge tiling by regular polygons containing at least one octagon is congruent
to the truncated square tiling $4.8.8$ ($t1002$). The truncated square tiling is
$1$-uniform; consequently no octagon occurs in any $k$-uniform tiling with $k \ge 2$.
\end{lemma}

\begin{proof}
Let $T$ contain an octagon. By O1 and O2, every vertex of $T$ incident to an octagon has
species $4.8.8$.

\emph{Step 1 (the vertex figure).} At a vertex $v$ of species $4.8.8$ the cyclic tile
order is (square, octagon, octagon). Hence: the two tiles adjacent to the square across
its two edges at $v$ are the two octagons; and each octagon has as cyclic neighbors at
$v$ the square on one side and the other octagon on the other.

\emph{Step 2 (squares adjacent to octagons).} Let a square $S$ share an edge $e$ with an
octagon. Both endpoints of $e$ are octagon-incident, hence of species $4.8.8$; by Step 1
the tiles across the two edges of $S$ adjacent to $e$ are octagons. Those octagons
contain the endpoints of the edge $e'$ of $S$ opposite to $e$, so those endpoints are
also of species $4.8.8$, and Step 1 there makes the tile across $e'$ an octagon too: all
four neighbors of $S$ are octagons, and all four vertices of $S$ are of species $4.8.8$.
(In particular no two squares share an edge: a $4.8.8$ vertex carries a single square.)

\emph{Step 3 (octagon coronas).} Let $O'$ be an octagon of $T$. Each of its vertices is
$4.8.8$, and by Step 1, at each vertex exactly one of the two edges of $O'$ borders a
square and the other an octagon. Consecutive edges of $O'$ share a vertex, so the
square/octagon assignment alternates around the eight edges; $8$ is even, so the
alternation is consistent, and the entire corona of $O'$ is determined by the assignment
at any single edge.

\emph{Step 4 (agreement induction).} A copy of the truncated square tiling containing a
given octagon is determined up to the corona parity of Step 3 (the point group of the
model at an octagon center is $D_4$, of index 2 in the $D_8$ of the octagon, and the two
cosets realize the two parities); let $M_0$ be the copy containing $O$ whose tile across
one fixed edge $e_0$ of $O$ matches $T$'s. Let $A$ be the set of tiles $t$ of $T$ such
that $t$ is a tile of $M_0$ \emph{and} every edge-neighbor of $t$ in $T$ equals the
corresponding tile of $M_0$. First, $O \in A$: its $T$-corona alternates (Step 3) with
the same pinning as $M_0$'s (choice of $M_0$), and an octagon or square sharing a
specified edge of $O$ from the outside is determined as a set by its shape and that
edge, so the corresponding tiles coincide. Now let $t \in A$ and let $t'$ be the tile of
$T$ across an edge $e$ of $t$; then $t' \in M_0$ (membership of $t$ in $A$), and we
check $t' \in A$. If $t'$ is an octagon, its $T$-corona alternates (Step 3) and is
pinned by its neighbor $t$, which is also its $M_0$-neighbor; determined-by-shape-and-edge
again makes each of its eight neighbors equal to $M_0$'s. If $t'$ is a square, then $t$
is an octagon, by the structure of $M_0$ rather than Step 2: $t \in A$ makes $t'$ the
$M_0$-neighbor of $t$ across $e$, and in the truncated square tiling $M_0$ a square is
edge-adjacent only to octagons, so a square neighbor $t'$ forces $t$ to be an octagon.
Now $t'$ is a square sharing edge $e$ with the octagon $t$, so Step 2 applies to $t'$ and
makes all four $T$-neighbors of $t'$ octagons; at each vertex $v$ of the shared edge $e$,
the third tile at $v$ is the octagon completing the $(4,8,8)$ figure, which shares a
specified edge with $t'$, and the fourth neighbor shares the opposite edge; each is
determined by shape and shared edge and so equals $M_0$'s tile. Hence $A$ is closed under edge-adjacency; the tile-adjacency graph
of a plane tiling is connected (join interior points of two tiles by a segment avoiding
vertices), so every tile of $T$ is a tile of $M_0$, and a tiling contained tile-wise in
another equals it: $T = M_0$.

\emph{1-uniformity.} In the truncated square tiling the point group at an octagon center
is $D_4$ and no vertex of the octagon lies on one of its mirrors (the mirrors pass
through edge midpoints), so $D_4$ acts freely, hence transitively, on the octagon's $8$
vertices; every vertex lies on an octagon, and all octagons are equivalent under
translations. So the symmetry group is vertex-transitive: exactly one vertex orbit.
\end{proof}

\begin{lemma}[O4: octagon-free implies palette $\{3,4,6,12\}$ and grid directions]
\label{lem:O4} \CLOSED\ (classical anchor: G\&S Table 2.1.1, pp.~59--61)\;
In an edge-to-edge tiling by regular polygons, only polygons with
$n \in \{3,4,6,8,12\}$ occur. Hence an octagon-free tiling uses
$\mathcal{P} = \{3,4,6,12\}$ only; all its interior angles are multiples of $30^\circ$,
and after a global rotation all edge directions lie in the $30^\circ$ grid.
\end{lemma}

\begin{proof}
The first sentence is classical, now pinned: G\&S enumerate the 17 species (21 types
counting cyclic order) of angle-feasible vertices in Table 2.1.1 and state on p.~59 that
for the six species marked with two asterisks there, namely $3.7.42$, $3.8.24$,
$3.9.18$, $3.10.15$, $4.5.20$, $5.5.10$, ``there is no edge-to-edge tiling by regular
polygons that includes even a single vertex of the species''; the surviving species use
only $n \in \{3,4,6,8,12\}$. Each of the six exclusions is an odd-polygon adjacency
argument of the O2 shape and can be reproved in an appendix if self-containedness is
preferred. Glue:
suppose the tiling is octagon-free, so all tiles lie in $\mathcal{P}$ and every interior
angle is $60, 90, 120$ or $150$ degrees, a multiple of $30$. Assign to each
\emph{edge-ray} (an edge directed away from one of its endpoints) its direction in
$\R/360$. At a fixed vertex, consecutive rays in the rotational order differ by interior
angles, so all rays at one vertex lie in a single coset of $30\Z$. The two rays of one
edge differ by $180 = 6 \cdot 30$, so the coset is shared along every edge. The graph
$(V(T), E(T))$ is connected (any two vertices are joined through the tile-adjacency chain
along a segment between them, avoiding vertices by general position), so all edge-rays of
$T$ lie in one coset $\alpha + 30\Z$; rotating by $-\alpha$ places every direction on the
grid. Finally, $t1002$ genuinely needs the $15^\circ$ grid: its vertex figure mixes $90$
and $135$ degree angles, so its rays occupy a coset of $45\Z$ meeting no single coset of
$30\Z$; this is why the truncated square tiling is outside the $\zeta_{12}$ model and
outside $\T_1$.
\end{proof}

\begin{proof}[Proof of Theorem~\ref{thm:octagon} from O1--O4]
A $k$-uniform tiling either contains an octagon, and is then $t1002$ with $k = 1$ (O3),
or contains none and lies in $\T_k$ by definition. That $t1002$ is $1$-uniform is part of
the phase close of O3 (its group $p4m$ acts transitively on vertices).
\end{proof}

%======================================================================
\section{The combinatorial model}\label{sec:model}

This section carries the definitional groundwork (the risk-4 item) and is now closed:
the stub-system axioms, the flag calculus of a tiling, the dictionary between the two,
and the quotient construction, each with proof.

\subsection{Stub systems}

\begin{definition}[stub system]\label{def:stub}
A \emph{stub system} is a tuple $(X, \rho, \mu, \gamma, c)$: a finite set $X$ of
\emph{stubs}; a bijection $\rho : X \to X$ (\texttt{rneig}; $\lambda := \rho^{-1}$ is
\texttt{lneig}); an involution $\mu : X \to X$ (\texttt{mirro}) with
$\mu\rho\mu = \rho^{-1}$; a partial map $\gamma : X \rightharpoonup X$ (\texttt{glue})
with $\gamma^2 = \mathrm{id}$ on its domain (fixed points allowed: self-gluings), domain
$\mu$-invariant and $\gamma\mu = \mu\gamma$ on the domain; and a coloring
$c : X \to \mathcal{C}$ (\texttt{lvert}; for the regular palette
$\mathcal{C} = \mathcal{P}$) satisfying $c \circ \mu = c \circ \rho$. A \emph{vertex} of
$X$ is an orbit of $\langle \rho, \mu \rangle$. $X$ is \emph{closed} iff $\gamma$ is
total, \emph{connected} iff $\langle \rho, \mu, \gamma \rangle$ acts transitively.
Morphisms are maps commuting with $\rho, \mu, \gamma$ (where defined) and preserving $c$;
isomorphisms accordingly.
\end{definition}

\begin{remark}[parity is an axiom consequence]\label{rem:parity}
In any stub system, $\gamma(x)$ is $\mu$-fixed iff $x$ is: $\mu\gamma x = \gamma\mu x$,
and $\gamma$ is injective on its domain, so $\mu\gamma x = \gamma x \iff \mu x = x$.
This is the abstract half of the mirror-parity rule (L2); no geometric input is needed.
\end{remark}

\begin{definition}[face walk, validity]\label{def:valid}
The \emph{face walk} is the partial map $\omega := \rho \circ \gamma$. A stub $x$ with
$n = c(x)$ satisfies:
\emph{condition 1} if $c$ is constant along the $\omega$-walk through $x$;
\emph{condition 2a} if, when the walk through $x$ is a cycle of minimal period $\ell$,
$\ell \mid n$;
\emph{condition 2b} if, when the walk is not a cycle, the maximal $\omega$-chain through
$x$ has length $\le n$ (length counts corners visited).
A partial stub system is \emph{valid} if all three hold at every stub where applicable.
\end{definition}

\begin{definition}[alphabet labeling]\label{def:overalpha}
A stub system is \emph{over the alphabet} $\A$ (Section~\ref{sec:alphabet}) if each of
its vertices, as a substructure $(\text{orbit}; \rho, \mu, c)$ with $\gamma$ forgotten,
is isomorphic to an entry of $\A$. By Lemma A6 the entry is then unique, so vertices are
canonically labeled.
\end{definition}

\begin{definition}[congruence, core, refinement]\label{def:congruence}
A \emph{congruence} on a stub system is an equivalence relation $\theta$ on $X$,
compatible with $\rho, \mu, \gamma$ ($x \,\theta\, y \Rightarrow f(x)\,\theta\,f(y)$,
with $\gamma$-definedness constant on classes) and refining the fibers of $c$. The
quotient $X/\theta$ is again a stub system. $X$ is a \emph{core} iff its only congruence
is equality. The \emph{refinement partition} $W(X)$ is the coarsest partition of $X$
refining the $c$-fibers and stable under $\rho, \lambda, \mu, \gamma$; it is what
\texttt{simplify} computes (Lemma R1).
\end{definition}

\subsection{Flags of a tiling and the dictionary}

\begin{definition}[flags]\label{def:flags}
A \emph{flag} of a tiling $T$ is an incident triple $x = (v, e, f)$, vertex--edge--face.
On $\Flags(T)$ define the three involutions
$\sigma_0$ (replace $v$ by the other endpoint of $e$),
$\sigma_1$ (replace $e$ by the other edge of $f$ at $v$),
$\sigma_2$ (replace $f$ by the other face at $e$).
In an edge-to-edge tiling each is well defined and fixed-point free, and
$\sigma_0\sigma_2 = \sigma_2\sigma_0$ (they alter disjoint components). The
\emph{dictionary} is
\[
\rho := \sigma_2\sigma_1, \qquad
\mu := \sigma_2, \qquad
\gamma := \sigma_2\sigma_0, \qquad
c(x) := \bigl|\,\text{face of } \sigma_2 x\,\bigr|
\]
(the color of a flag is the size of the \emph{opposite} face across its edge). Every
isometry in $G(T)$ commutes with $\sigma_0, \sigma_1, \sigma_2$ (they are
incidence-defined) and preserves face sizes, hence commutes with the dictionary;
orientation plays no role at the flag level, which is why flags rather than oriented
darts are the right carrier and why Convention~\ref{conv:chiral} is native to the model.
\end{definition}

\begin{remark}[correction to the phase-1 skeleton]\label{rem:gammafix}
Phase 1 declared $\gamma := \sigma_0$. That is wrong: with $\gamma = \sigma_0$ the face
walk $\omega = \rho\gamma$ changes faces at every step and condition 1 would fail in
$\Flags(T)$ itself. The computation in D1a below forces $\gamma = \sigma_2\sigma_0$: when
an edge is crossed, the local rotational sense reverses, and the extra $\sigma_2$ is that
reversal. The error was caught by carrying out the walk computation rather than trusting
the analogy with Delaney--Dress conventions; no lemma statement was affected.
\end{remark}

\begin{lemma}[D1a: flag calculus]\label{lem:D1a} \CLOSED\;
Let $T$ be a tiling. Then $(\Flags(T); \rho, \mu, \gamma, c)$ satisfies the stub-system
axioms with $\gamma$ total, is connected, and is valid; moreover the $\omega$-orbits are
in bijection with the faces of $T$: the walk through a stub of color $n$ is a cycle of
length exactly $n$ visiting the $n$ corners of one face in rotational order.
\end{lemma}

\begin{proof}
\emph{Axioms.} $\mu^2 = \sigma_2^2 = \mathrm{id}$. $\gamma^2 =
\sigma_2\sigma_0\sigma_2\sigma_0 = \mathrm{id}$ by commutation. $\gamma\mu =
\sigma_2\sigma_0\sigma_2 = \sigma_2\sigma_2\sigma_0\ldots$ more directly $\gamma\mu =
\sigma_2\sigma_0\sigma_2$ and $\mu\gamma = \sigma_2\sigma_2\sigma_0 = \sigma_0$; using
commutation $\sigma_2\sigma_0\sigma_2 = \sigma_0\sigma_2\sigma_2 = \sigma_0$, so
$\gamma\mu = \mu\gamma$. $\mu\rho\mu = \sigma_2(\sigma_2\sigma_1)\sigma_2 =
\sigma_1\sigma_2 = (\sigma_2\sigma_1)^{-1} = \rho^{-1}$.

\emph{Local coordinates.} Orient the plane once and for all. At a vertex $v$ of degree
$d$, list the incident edges $e_0, \dots, e_{d-1}$ in counterclockwise order (indices mod
$d$) and let $f_i$ be the face between $e_i$ and $e_{i+1}$. The $2d$ flags at $v$ are
$(v, e_i, f_i)$ and $(v, e_i, f_{i-1})$; write them $(i,0)$ and $(i,1)$. Direct
computation: $\sigma_1(v,e_i,f_i) = (v,e_{i+1},f_i)$ and then $\sigma_2$ gives
$(v,e_{i+1},f_{i+1})$, so $\rho(i,0) = (i{+}1,0)$; similarly $\rho(i,1) = (i{-}1,1)$;
and $\mu(i,b) = (i,1{-}b)$. Colors: $c(i,0) = |f_{i-1}|$ and $c(i,1) = |f_i|$, and one
checks $c\mu = c\rho$ on both chirality types. (These formulas coincide with the
generator's doubled dart structure; see A2.)

\emph{Crossing an edge.} Let $e$ carry index $a$ at $v$ and index $b$ at the other
endpoint $w$. Counterclockwise order at $v$ appears clockwise from $w$, which gives the
face matching $f^v_a = f^w_{b-1}$ and $f^v_{a-1} = f^w_b$. Hence
$\sigma_0(v,e_a,f^v_a) = (w,e_a,f^w_{b-1})$, a type-1 flag at $w$, and
$\gamma(a,0)_v = \sigma_2\sigma_0(v,e_a,f^v_a) = (w, e_a, f^w_b) = (b,0)_w$: $\gamma$
preserves chirality type relative to the global orientation.

\emph{Face walk.} Take the corner $(\lambda x, x)$ with $x = (i{+}1,0)$ at $v$; its color
is $c(x) = |f^v_i|$. Then $\gamma(x) = (j,0)$ at $w$ (the other end of $e_{i+1}$, with
$f^v_i = f^w_j$ by the crossing rule), and $\omega(x) = \rho\gamma(x) = (j{+}1, 0)$ at
$w$, whose color is $|f^w_j| = |f^v_i|$: the walk has advanced one corner around the same
geometric face $f := f^v_i$ and kept its color. Iterating, the walk visits the corners of
$f$ in rotational order; since $f$ is an embedded $n$-gon with $n$ distinct corners, the
walk first returns after exactly $n$ steps. Conversely every corner of every face occurs
in such a walk, giving the bijection. Conditions 1, 2a hold with $\ell = n$; there are no
open chains ($\gamma$ total). The same computation on type-1 flags gives the reverse
walks; nothing changes.

\emph{Connectivity.} Flags at one vertex are connected by $\rho, \mu$. Any two vertices
are joined by a path of edges (O4's connectivity argument), and $\gamma$ crosses each
edge. Hence $\langle \rho, \mu, \gamma\rangle$ is transitive.
\end{proof}

\begin{lemma}[B0: flag rigidity and discreteness]\label{lem:B0} \CLOSED\;
A symmetry of $T$ fixing a flag is the identity; $G(T)$ acts freely on $\Flags(T)$; the
vertex set of $T$ is uniformly discrete and $G(T)$ is a discrete group.
\end{lemma}

\begin{proof}
Let $g$ fix $(v,e,f)$: $g(v) = v$ and $g(e) = e$, so $g$ fixes both endpoints of $e$,
hence fixes the line of $e$ pointwise; an isometry fixing a line pointwise is the
identity or the reflection in that line, and the reflection swaps the two faces at $e$,
contradicting $g(f) = f$. So $g = \mathrm{id}$, and the action on flags is free.
Uniform discreteness: two distinct vertices have distance $\ge 1$. Indeed if
$0 < d(v,w) < 1$, then $w$ lies within distance $1$ of $v$; every point at distance
$< \tfrac12$ from $v$ lies in a tile incident to $v$ (each unit-edge regular tile
contains the disk sector of radius $\ge \sqrt3/2$ at each of its corners), and for
$d(v,w) < 1$ one checks $w$ lies on the boundary of a tile $P$ incident to $v$; by the
edge-to-edge property a vertex of the tiling on $\partial P$ is a corner of $P$, and
distinct corners of a unit-edge regular polygon are at distance $\ge 1$ from each other,
a contradiction. Discreteness of $G(T)$: if $g_m \to \mathrm{id}$, then for large $m$,
$g_m$ moves a fixed flag's vertex, edge midpoint and face barycenter by less than the
minimal separation of those data, forcing $g_m$ to fix that flag, so $g_m =
\mathrm{id}$.
\end{proof}

\begin{lemma}[D1b: quotients]\label{lem:D1b} \CLOSED\;
Let $H \le G(T)$ have finitely many orbits on $\Flags(T)$. Then
$T/H := (\Flags(T)/H, \text{induced operations})$ is a finite closed connected stub
system; its vertices correspond to $H$-orbits of $V(T)$. Moreover, for a stub
$\bar{x} = Hx$ with $x = (v,e,f)$:
\begin{itemize}
\item[(i)] $\mu\bar x = \bar x$ iff some $h \in H$ is the reflection in the line of $e$
  (a mirror along the edge, through both endpoints);
\item[(ii)] $\gamma\bar x = \bar x$ iff some $h \in H$ is the rotation by $\pi$ about the
  midpoint of $e$ (the $(a)$-gluing);
\item[(iii)] $\gamma\bar x = \mu\bar x$ iff some $h \in H$ is the reflection in the
  perpendicular bisector of $e$ (the $[a]$-gluing).
\end{itemize}
Finally $T \mapsto T/H$ is functorial: an isometry $g$ with $gHg^{-1} = H'$ induces an
isomorphism $T/H \cong (gT)/H'$; in particular mirror-image tilings have isomorphic
canonical quotients (chirality merge at the model level).
\end{lemma}

\begin{proof}
$H$ commutes with the dictionary (Definition~\ref{def:flags}), so the operations descend;
the axioms and closedness are inherited, connectivity passes to quotients, and finiteness
is the orbit hypothesis. Vertices: $\langle\rho,\mu\rangle$-orbits in the quotient are
images of flag sets at single vertices, and two vertices' flag sets meet the same orbit
iff the vertices are $H$-equivalent. For (i): $\mu\bar x = \bar x$ means $hx = \sigma_2 x$
for some $h \in H$; such $h$ fixes $v$ and $e$, hence the line of $e$ pointwise (as in
B0), and is not the identity (the faces differ), so it is the reflection in that line;
conversely that reflection realizes $hx = \sigma_2x$. For (ii): $hx = \sigma_2\sigma_0 x$
maps $(v,e,f)$ to $(w,e,f')$ with $w \ne v$, $f' \ne f$: $h$ preserves $e$, swaps its
endpoints and swaps its sides, so $h$ is the half-turn about the midpoint; conversely the
half-turn does this. For (iii): $hx = \mu\gamma\bar x$-wise, $hx = \sigma_0x$ maps
$(v,e,f)$ to $(w,e,f)$: $h$ preserves $e$ and $f$, swaps endpoints, so it fixes the
midpoint and the perpendicular direction: the reflection in the perpendicular bisector.
Functoriality: $g$ maps flags of $T$ to flags of $gT$ commuting with everything, and
descends to the orbit spaces. For the chirality statement take $H = G(T)$,
$g$ any reflection: $(gT)/G(gT) \cong T/G(T)$.
\end{proof}

\begin{remark}[presentation vs.\ structure: starred labels]\label{rem:presentation}
The implementation presents each vertex orbit by choosing one ``orientation sheet'':
stubs of the chosen sheet are unstarred, their $\mu$-images starred. The choice is
arbitrary per orbit, and nothing in the structure depends on it; in particular
$\gamma$ may connect an unstarred stub to a starred one (the $[a\,b]$ Conway form)
exactly when the two vertex orbits' chosen sheets disagree across the edge, since
geometric $\gamma$ preserves the \emph{global} chirality type (D1a, crossing rule) while
the sheets are local conventions. The four Conway bracket forms are therefore
presentation data; the invariant content is D1b (i)--(iii). This is also why the pruner's
isomorphism test, which never reads labels, is the right notion of sameness.
\end{remark}

\begin{lemma}[L1: necessity and heredity]\label{lem:L1} \CLOSED\;
(i) For any tiling $T$ and $H \le G(T)$ with finitely many flag orbits, $T/H$ is valid:
every $\omega$-cycle has constant color $n$ and length $\ell = n/r$ dividing $n$, where
$r$ is the order of the rotation part of the $H$-stabilizer of the corresponding face.
(ii) Heredity: if $M$ is closed and valid and $P$ is a partial substructure of $M$ (a
subset of vertices carrying a subset of the gluings), then $P$ is valid.
\end{lemma}

\begin{proof}
(i) Fix a stub $\bar x$ of $T/H$ and lift to $x \in \Flags(T)$; by D1a the planar walk
through $x$ runs around a face $f$ with $|f| = n = c(x)$, visiting flags
$x, \omega x, \dots, \omega^{n-1}x$, all of color $n$; colors descend, giving condition 1.
The quotient walk closes at the minimal $\ell$ with $\omega^\ell x \in Hx$, say
$h x = \omega^\ell x$ with $h \in H$; since $H$ commutes with $\omega$, $h$ maps the
walk orbit to itself and stabilizes $f$. \emph{Claim: such $h$ is orientation-preserving.}
Define for a flag $y = (u, e, f)$ the local orientation $\mathrm{or}(y) \in \{\pm 1\}$ as
the orientation of the frame (direction of $e$ away from $u$, direction from $e$ into
$f$). Each $\sigma_i$ reverses $\mathrm{or}$: the defining data of $y$ and $\sigma_i y$
are exchanged by an explicit reflection (in the angle bisector of the corner for
$\sigma_1$, the perpendicular bisector of $e$ for $\sigma_0$, the line of $e$ for
$\sigma_2$), and $\mathrm{or}$ is covariant under isometries with the sign of their
determinant. Hence $\omega = \sigma_2\sigma_1\sigma_2\sigma_0$ preserves $\mathrm{or}$,
so all flags in one $\omega$-orbit share their value; from
$\mathrm{or}(hx) = \det(h)\,\mathrm{or}(x)$ and $hx$ in the orbit of $x$ we get
$\det(h) = +1$, proving the claim. So the return elements form the cyclic rotation group
$C_r = \Stab_H(f) \cap SO$, acting on the $n$ corners of $f$ freely with orbits of size
$r$; the minimal return is $\ell = n/r$, which divides $n$: condition 2a. Open chains do
not occur ($\gamma$ total). Reflections in $\Stab_H(f)$, if any, do not shorten the
cycle; they surface as coincidences between the cycle and its $\mu$-image.
(ii) Let $\iota : P \to M$ be the inclusion-style morphism ($\rho,\mu,c$ preserved,
$\gamma_P$ a restriction of $\gamma_M$). Walks in $P$ are walks in $M$ while defined and
$\iota$ is injective, so: a $P$-cycle is an $M$-cycle with the same minimal period
(deterministic walks return in $P$ iff they return in $M$ at the same step), giving 2a;
a maximal open $P$-chain visits distinct stubs injectively inside an $M$-cycle of length
$\ell \le n$, giving 2b; colors are inherited, giving 1. Monotonicity for the search:
gluings are only ever added along a DFS branch, walks only extend, so a violation of
1/2a/2b in a configuration persists in all its descendants; equivalently no ancestor of a
valid configuration violates the conditions.
\end{proof}

\begin{lemma}[L2: mirror parity]\label{lem:L2} \CLOSED\;
In any stub system, $\gamma$ maps $\mu$-fixed stubs to $\mu$-fixed stubs
(Remark~\ref{rem:parity}); in $T/H$ a stub is $\mu$-fixed iff a mirror of $H$ runs along
its edge (D1b(i)), and the mirror line of an edge contains both endpoints, so the two
end-stubs of such an edge are both $\mu$-fixed. Consequently the solver's parity filter
(glue self-mirrored only to self-mirrored) excludes no gluing that occurs in any $T/H$,
and is forced for every closed valid system.
\end{lemma}

\begin{proof}
All three assertions are contained in Remark~\ref{rem:parity} and D1b(i); the last
sentence restates them for the search.
\end{proof}

\begin{remark}[kinship with Delaney--Dress symbols]
$\Flags(T)/G(T)$ with $(\sigma_0, \sigma_1, \sigma_2)$ is the underlying set of the
Delaney--Dress symbol of $T$; the stub system repackages the operators as
$(\rho, \mu, \gamma)$ and carries sizes as colors. The correspondence theorems (C, B) are
vertex-centered analogues of the D-symbol realization theorems (Dress 1987, Dress--Huson
1987), cited as context and pattern only; the self-contained 2D proofs are shorter than a
formal bridge between the conventions would be.
\end{remark}

%======================================================================
\section{Obligation 1: the alphabet}\label{sec:alphabet}

The alphabet $\A$ (\texttt{mainlist}) is the finite set of vertex letters: pairs
(configuration, site-symmetry variant), encoded as stub structures with the attachment
data \texttt{ferkval}/\texttt{reps}. The generator \texttt{alphabets/gen\_alphabet.py}
constructs $\A$ from first principles and gates it against the 44 legacy entries.

\begin{lemma}[A1: configuration enumeration]\label{lem:A1} \CLOSED\;
The cyclic words over $\mathcal{P} = \{3,4,6,12\}$ with interior angles summing to
$360^\circ$, up to rotation and reflection, are exactly the following 14:
\[
\begin{array}{llll}
(3,12,12), & (4,6,12), & (6,6,6), & (4,4,4,4),\\
(3,3,6,6), & (3,6,3,6), & (3,4,4,6), & (3,4,6,4),\\
(3,3,4,12), & (3,4,3,12), & (3,3,3,3,6), & (3,3,3,4,4),\\
(3,3,4,3,4), & (3,3,3,3,3,3), &&
\end{array}
\]
of which exactly $(4,6,12)$, $(3,3,4,12)$ and $(3,4,4,6)$ are chiral (not fixed by any
reflection of the cyclic word).
\end{lemma}

\begin{proof}
With $m$ polygons at a vertex, the angle equation reads
$\sum_i (1 - 2/n_i) = 2$, i.e.\ $\sum_i 1/n_i = (m-2)/2$ with
$1/n_i \in \{1/3, 1/4, 1/6, 1/12\}$. Since each term is at most $1/3$, $m/3 \ge (m-2)/2$
forces $m \le 6$; since each term is positive and at most... for $m \le 2$ no solution
exists. Case $m=3$ ($\sum = 1/2$): with a $1/3$: remaining $1/6$ from two terms:
$1/12 + 1/12$ only, giving $\{3,12,12\}$; without: $1/4+1/6+1/12 = 1/2$ giving
$\{4,6,12\}$, and $1/6 \cdot 3 = 1/2$ giving $\{6,6,6\}$; all other combinations miss.
Case $m=4$ ($\sum = 1$), by the number $t$ of triangles: $t = 2$: remaining $1/3$ from
two terms: $1/4 + 1/12$ ($\{3,3,4,12\}$) or $1/6+1/6$ ($\{3,3,6,6\}$); $t = 1$: remaining
$2/3$ from three terms of $\{1/4,1/6,1/12\}$: only $1/4+1/4+1/6$ ($\{3,4,4,6\}$);
$t = 0$: only $1/4 \cdot 4$ ($\{4,4,4,4\}$); $t \ge 3$ overshoots. Case $m=5$
($\sum = 3/2$): $t = 4$: fifth term $1/6$ ($\{3,3,3,3,6\}$); $t = 3$: two terms summing
$1/2$: only $1/4+1/4$ ($\{3,3,3,4,4\}$); $t \le 2$: maximum $2/3 + 3/4 < 3/2$. Case
$m=6$: all six terms must be $1/3$ ($\{3^6\}$). Cyclic arrangements up to the dihedral
action: 3-element multisets have one arrangement each; $\{3,3,6,6\}$: equal pairs
adjacent or alternating, $2$; $\{3,4,4,6\}$: the squares adjacent or separated, $2$;
$\{3,3,4,12\}$: triangles adjacent or separated, $2$ (adjacent: the two rotation classes
$(3,3,4,12)$ and $(3,3,12,4)$ merge under reflection, leaving one chiral dihedral class);
$\{4,4,4,4\}, \{3,3,3,3,6\}, \{3^6\}$: one each; $\{3,3,3,4,4\}$: the lone triangle
flanked by squares or not, i.e.\ $(3,3,3,4,4)$ and $(3,3,4,3,4)$, $2$. Total
$3 + 7 + 3 + 1 = 14$. Chirality: a cyclic word is achiral iff some rotation of its
reversal equals it; checking each, exactly the three listed fail, e.g.\ the rotations of
$(3,3,4,12)$ are $(3,3,4,12), (3,4,12,3), (4,12,3,3), (12,3,3,4)$ and its reversal is
$(3,3,12,4)$ up to rotation, not among them. (Consistency with the classical count: of
G\&S's 21 types over all polygon sizes, 15 occur in tilings; removing the type $4.8.8$
leaves the 14 above.)
\end{proof}

\begin{lemma}[A2: variants are subgroup conjugacy classes]\label{lem:A2} \CLOSED\;
Fix a configuration $w$ of length $m$ and let $\Sym(w) \le D_m$ be the stabilizer of the
cyclic word under the dihedral action on positions. For a tiling $T$ and a vertex $v$
with configuration $w$:
\begin{itemize}
\item[(i)] $\Stab_{G(T)}(v)$ acts faithfully on the flags at $v$, and under the local
  coordinates of D1a this action is by word symmetries: rotations act as
  $(i,b) \mapsto (i+t, b)$ and reflections as $(i,b) \mapsto (a-i, 1-b)$; the image is a
  subgroup $U \le \Sym(w)$.
\item[(ii)] The vertex substructure of $T/H$ at $v$ (for any $H$, with
  $U = $ image of $\Stab_H(v)$) is the orbit space of the doubled dart structure of $w$
  under $U$, i.e.\ the generator's fold.
\item[(iii)] Conjugate subgroups of $\Sym(w)$ yield isomorphic folds; hence the map
  \{variants\} $\to$ \{conjugacy classes of subgroups of $\Sym(w)$\} is a well-defined
  bijection onto, and no realizable vertex letter is missing from an alphabet that lists
  one fold per conjugacy class.
\item[(iv)] (angle-weight clause, consumed by the C-block) With $a(n) := 6 - 12/n$, every
  $\rho$-cycle $C$ of the fold of $w$ by $U$ has angle-weight sum
  $s(C) := \sum_{x \in C} a(c(x)) = 12/r$, where $r$ is the order of the rotation
  subgroup of $U$; in particular $s(C)$ divides $12$.
\end{itemize}
\end{lemma}

\begin{proof}
(i) An element $h$ of the stabilizer fixes $v$, hence is a rotation about $v$ or a
reflection through $v$; it permutes the incident edges preserving their counterclockwise
cyclic order up to sense, and preserves the sizes of the faces between them, so its
action on indices is an element of $D_m$ stabilizing the word $w$; the formulas for the
two chirality types follow from D1a's coordinates ($h$ orientation-preserving fixes each
sheet, orientation-reversing swaps them). Faithfulness: $h$ acting trivially on flags at
$v$ fixes a flag, hence $h = \mathrm{id}$ by B0. (ii) The flags at $v$ \emph{are} the
doubled dart set $\{(i,b)\}$ with the operations computed in D1a, which coincide with the
generator's $\texttt{rneig}(i,0)=(i{+}1,0)$, $\texttt{rneig}(i,1)=(i{-}1,1)$,
$\texttt{mirro}(i,b)=(i,1{-}b)$, $\texttt{cls}(i,0)=c[i{-}1]$, $\texttt{cls}(i,1)=c[i]$;
the quotient by $\Stab_H(v)$ is by definition the fold by $U$. (iii) If
$U' = sUs^{-1}$ with $s \in \Sym(w)$, then $s$ induces an isomorphism of folds
$\text{darts}/U \to \text{darts}/U'$, $Ux \mapsto U'sx$, which respects the operations
because $s$ does. Surjectivity of the classification is (i): every stabilizer that occurs
is a subgroup of $\Sym(w)$, so its conjugacy class is listed. As a consistency check
against the shipped table: $\Sym = 1$ for the three chiral words (1 variant each),
$\Sym = \Z_2$ for the seven reflection-only words (2 variants), $\Sym(3,6,3,6) = D_2$
(5 subgroups, all normal: F, A1, A2, R, S), $\Sym(6,6,6) = D_3$ (4 classes),
$\Sym(4,4,4,4) = D_4$ (8 classes), $\Sym(3^6) = D_6$ (10 classes; the two $D_3$ classes
are the a/b pair, and the single $D_2$ class explains why $(3^6)S2$ has one edge-type
axis and one triangle-type axis), for a total of
$3 + 14 + 5 + 4 + 8 + 10 = 44$, the legacy entry count.

(iv) Lift a $\rho$-cycle of the fold to the doubled dart structure: $\rho$ preserves the
two sheets there, so the lifted walk runs within one sheet, stepping through the word
positions in one rotational sense with corner weights $a(c_j)$, whose total over the
whole word is $12$ (the angle equation of A1, in $30^\circ$ units). The folded cycle
closes at the least $t$ for which the lifted dart returns modulo $U$; sheet-preserving
elements of $U$ are exactly its rotations, which form the cyclic group generated by the
shift by $m/r$, so $t = m/r$. Rotations preserve the color word, so the color word is
$(m/r)$-periodic and the partial weight sum over $m/r$ consecutive corners is $12/r$.
Sheet-1 cycles are the mirror images of sheet-0 cycles and carry the same weights.
\end{proof}

\begin{remark}[spurious variants are harmless]
The brief asks that ``none of the generated variants is spurious.'' That direction is not
load-bearing: a letter that never occurs in a valid closed core simply never appears in
output, and soundness of what does appear is Obligation 5's job, not the alphabet's. Only
missing letters (A2(iii)) and wrongly encoded letters (A3/A4) can break the theorem.
\end{remark}

\begin{lemma}[A3: table correctness]\label{lem:A3} \MACHINE\ (implemented, unaudited)\;
The generated (configuration $\times$ subgroup class) folds coincide, arrays and
\texttt{ferkval} included, with the 44 legacy \texttt{mainlist} entries, trying every
frame per fold; and the three consumed table copies (\texttt{solver\_tables.inc},
\texttt{pruner\_tables.inc}, the original hard-coded \texttt{mainlist}) agree.
\end{lemma}

\begin{lemma}[A4: structural identities]\label{lem:A4} \MACHINE\ (implemented, unaudited)\;
Every entry satisfies $\lambda\rho = \mathrm{id}$, $\mu^2 = \mathrm{id}$,
$\mu\rho = \lambda\mu$, $c\,\mu = c\,\rho$.
\end{lemma}

\begin{lemma}[A5: the ferk transversal]\label{lem:A5} \MACHINE\ (implemented, unaudited)\;
For every entry $D$: $|\Aut(D)| = \texttt{ferkval}(D)$, the action of $\Aut(D)$ on stubs
is free, and the generated \texttt{reps} list meets every $\Aut(D)$-orbit.
\end{lemma}

\begin{lemma}[A6: pairwise non-isomorphism --- new certificate]\label{lem:A6}
\MACHINE\ (\textbf{to be added})\;
The 44 entries are pairwise non-isomorphic as colored stub structures
$(\text{orbit}; \rho, \mu, c)$.
\end{lemma}

\begin{remark}[why A6 is needed, and that it was missing]
The pruner buckets by the vertex-letter signature string and only compares within a
bucket; the letter multiset must therefore be an isomorphism invariant of the underlying
structure, which is exactly A6. Without it an over-count would be possible in principle:
one tiling emitted under two letterings, landing in different buckets, kept twice. A6 is
finite ($44^2$ colored-structure isomorphism tests) but is \emph{not} among the
generator's current certificates; this note adds it as an obligation. Spot-check done by
hand: the candidate collision $(3,6,3,6)A1$ vs $(3,6,3,6)A2$ is refuted by the corner
colors at the $\mu$-symmetric corners.
\end{remark}

%======================================================================
\section{Obligation 2: local rules}\label{sec:local}

L1 and L2 are stated and proved in Section~\ref{sec:model} (Lemmas \ref{lem:L1},
\ref{lem:L2}), where they belong to the quotient analysis. What remains here is the
scope remark:

\begin{remark}[L3: sufficiency is the development theorem]
The brief's ``conditions jointly sufficient for a closed configuration to be a valid
combinatorial tiling'' has no separate combinatorial content: the honest sufficiency
statement is exactly the C-block (a closed valid rigid system develops to a genuine
tiling), and it is proved there. Listing it here would double-count an obligation.
\end{remark}

%======================================================================
\section{Obligation 3: search exhaustiveness}\label{sec:search}

The abstract DFS. State: a connected partial stub system $P$ over $\A$ with an ordered
vertex list (letters indexed into $\A$). Root: a single fresh letter, no gluings. Step:
pick a free stub $s$ by any deterministic rule (the implementation uses
\texttt{mincycle}, most-constrained-first, and normalizes $s$ to the unstarred partner;
the proofs below are indifferent to the rule); enumerate as children:
(a) for every free stub $t$ with the same $\mu$-parity as $s$, the configuration
$P + \{s \leftrightarrow t\}$, with the mirror pair $\mu s \leftrightarrow \mu t$ glued
alongside when $s \ne \mu s$; the cases $t = s$ (self-gluing) and $t = \mu s$ are
included;
(b) if the vertex budget allows, for every letter $g \ge$ the root letter and every
$r \in \texttt{reps}(g)$ of matching parity, the configuration
$P + (\text{fresh vertex of letter } g) + \{s \leftrightarrow r\}$ (mirror pair
alongside).
Children failing conditions 1/2a/2b are discarded; closed children go to the rigidity
filter; open children recurse.

\begin{lemma}[T1: termination]\label{lem:T1} \CLOSED\;
The DFS tree is finite and the run halts.
\end{lemma}

\begin{proof}
A configuration carries at most \texttt{maxnum} vertices of at most 12 stubs each (the
largest letter, $(3^6)$F, has 12), so at most $12\,\texttt{maxnum}$ stubs. Every child
adds at least one gluing and none is ever removed along a branch, so branch length is at
most the stub count. Branching at a node is bounded by the number of free stubs plus
$|\A| \cdot \max_g |\texttt{reps}(g)| < \infty$. A finitely-branching tree with bounded
branch length is finite; the root loop over the 44 letters is finite.
\end{proof}

\begin{lemma}[S1: guided descent --- the no-drop core]\label{lem:S1} \CLOSED\;
Let $M$ be a finite connected closed valid stub system over $\A$ with
$m \le \texttt{maxnum}$ vertices whose minimal letter is $t^\ast$. Then the DFS rooted at
$t^\ast$ visits at least one closed configuration isomorphic to $M$.
\end{lemma}

\begin{proof}
Say a configuration $P$ \emph{follows} $M$ via $\psi$ if $\psi : X_P \to X_M$ is
injective, preserves $\rho, \mu, c$, and carries each gluing of $P$ to a gluing of $M$
(so $\gamma_M \psi = \psi \gamma_P$ on $\mathrm{dom}\,\gamma_P$). Since $\psi$ preserves
$\rho$ and $\mu$ and letters are single $\langle\rho,\mu\rangle$-orbits, $\psi$ maps each
vertex of $P$ bijectively onto a vertex of $M$ of the same letter, and distinct vertices
to distinct vertices.

\emph{Base.} $M$ has a vertex of letter $t^\ast$; the root configuration (one fresh
$t^\ast$, no gluings) follows $M$ via any letter isomorphism onto it, which exists by
Definition~\ref{def:overalpha}.

\emph{Inductive step.} Suppose the DFS visits a non-closed $P$ that follows $M$ via
$\psi$; let $s$ be the free stub the rule picks. $M$ is closed, so
$y := \gamma_M(\psi s)$ exists. By Remark~\ref{rem:parity} applied in $M$, $y$ is
$\mu$-fixed iff $\psi s$ is, iff $s$ is ($\psi$ preserves $\mu$): parities match.

Case (a): $y = \psi(t)$ for some $t \in X_P$. Then $t$ is free in $P$: otherwise
$\psi\gamma_P(t) = \gamma_M(y) = \psi s$ with $\psi$ injective would give
$\gamma_P(t) = s$, contradicting $s$ free. The child $P' := P + \{s \leftrightarrow t\}$
(mirror pair alongside if $s \ne \mu s$) is among the enumerated children, and $\psi$
still follows: the new gluings map to $\{\psi s \leftrightarrow y\}$ and, for the mirror
pair, to $\{\mu\psi s \leftrightarrow \mu y\}$, which is a gluing of $M$ because
$\gamma_M$ commutes with $\mu$.

Case (b): $y$'s vertex is not in the image of $\psi$. Its letter $g$ satisfies
$g \ge t^\ast$ because $t^\ast$ is minimal among $M$'s letters, and $P$ has fewer than
$m \le \texttt{maxnum}$ vertices (its vertices inject into $M$'s and miss $y$'s), so the
fresh-vertex branch is enabled and tries letter $g$. The set of letter isomorphisms
$D_g \to (\text{$y$'s vertex in } M)$ is a nonempty torsor under $\Aut(D_g)$, so the
stubs $r$ with some isomorphism sending $r \mapsto y$ form exactly one
$\Aut(D_g)$-orbit; by A5 the \texttt{reps} list meets it, say at $r_0$ with isomorphism
$\psi_0$, $\psi_0(r_0) = y$. Parity of $r_0$ matches ($\psi_0$ preserves $\mu$-fixedness
and $y$'s parity matches $s$'s). The child
$P' := P + (\text{fresh } g) + \{s \leftrightarrow r_0\}$ is enumerated, and
$\psi' := \psi \cup \psi_0$ follows $M$ ($\psi_0$'s image vertex is new, so injectivity
holds; the gluing maps correctly as in case (a)).

\emph{No prune.} A configuration following $M$ is isomorphic (via $\psi$, onto its image)
to a partial substructure of $M$, hence valid by L1(ii); so \texttt{checkpart} never
discards a child produced above.

\emph{Closure.} Each step adds a gluing of $M$'s finite gluing set, so after finitely
many steps the guided branch reaches a closed $P^\ast$ following $M$ via some
$\psi^\ast$. Its image is closed under $\rho, \mu$ (unions of vertices) and under
$\gamma_M$ ($P^\ast$ closed and $\gamma_M\psi^\ast = \psi^\ast\gamma_{P^\ast}$), hence
equals $X_M$ by connectedness of $M$; an injective surjective morphism is an
isomorphism. Since the DFS explores every child of every visited node and the root is
visited, every node of the guided branch is visited, in particular $P^\ast$.
\end{proof}

\begin{lemma}[S2: minimal-letter rooting; sharding]\label{lem:S2} \CLOSED\;
The run roots one DFS at every letter $i \in \A$. The search rooted at $i$ visits exactly
(up to isomorphism) the closed valid systems over $\A$ whose minimal letter is $i$; in
particular every $M$ as in S1 is reached from its minimal letter, and the root sets of
distinct roots are disjoint. Sharding the root loop (\texttt{EU\_SHARD\_N/W}) partitions
the roots; the union of the shards' outputs equals the sequential run's output.
\end{lemma}

\begin{proof}
Containment one way: the fresh-vertex branch only adds letters
$g \ge \mathrm{vertype}[0] = i$ and the root letter $i$ is never removed, so every
configuration in the tree rooted at $i$ has letter multiset with minimum exactly $i$.
Containment the other way is S1 applied with root $t^\ast(M) = i$. Disjointness: the
minimal letter is an isomorphism invariant. Sharding assigns each root to exactly one
worker and workers share no state, so the multiset union of emissions over workers equals
the sequential emission multiset; the pruner's output is invariant under this
repartition because its dedup is by isomorphism class within $k$ (P1/P2), which is
insensitive to emission order across families. (The last step's implementation
correctness, byte-identity of the merged catalog, is an audit item; the mathematical
statement here is set equality of the pruned classes.)
\end{proof}

\begin{remark}[correcting the brief's description]
The brief describes a ``nondecreasing-vertex-type tie-break.'' The code constrains new
letters only by $g \ge \mathrm{vertype}[0]$ (min-root invariant,
\texttt{eu\_solver.cpp:680}); letters may otherwise arrive in any order. This is weaker
than nondecreasing order and makes S1/S2 easier, since the descent never needs to reorder
$M$'s vertices.
\end{remark}

\begin{lemma}[S3: emission gates]\label{lem:S3} \CLOSED\ (abstract) + audit-facing\;
In the abstract pipeline, every closed configuration visited with
$1 \le \#\text{vertices} \le \texttt{maxnum}$ that passes the rigidity filter is emitted.
On the regular palette no budget, cap, or filter other than conditions 1/2a/2b and the
rigidity filter gates emission.
\end{lemma}

\begin{proof}
Abstractly this is the definition of $\mathrm{Solve}$
(Section~\ref{sec:scope}); there is nothing to prove. The content is entirely
implementation-side and is carried by audit hooks 4 and 7: \texttt{seen} $= 0$ emits
every $k \in [1, \texttt{maxnum}]$; the noncounting budget is inert on the regular
palette (\texttt{has\_noncounting} false); \texttt{EU\_STREAM} changes the sink only;
\texttt{EU\_KONLY} is a user-selected per-$k$ restriction applied before dedup and does
not affect the selected $k$.
\end{proof}

\begin{lemma}[S4: the picked stub exists; the local tests are the conditions]
\label{lem:S4} \CLOSED\ (abstract) + audit-facing\;
In a non-closed configuration every free stub lies on an open $\omega$-walk, so a
most-constrained free stub is well defined; and the abstract DFS prunes exactly on
conditions 1/2a/2b. The choice rule affects traversal order only: S1's proof is
indifferent to which free stub is picked.
\end{lemma}

\begin{proof}
A free stub $x$ has $\gamma(x)$ undefined, so the maximal $\omega$-chain ending at
$\rho$-predecessors of $x$ is open, and $x$ bounds it; hence open walks exist iff free
stubs exist. The second and third assertions restate the abstract pipeline definition
and S1. That \texttt{checkpart} decides exactly 1/2a/2b (walking each cycle from every
starting stub, so every maximal open chain is checked at full length from its earliest
stub) and that \texttt{mincycle} returns a free stub are implementation facts: audit
hook 5.
\end{proof}

%======================================================================
\section{Obligation 4: rigidity filter and dedup}\label{sec:rigidity}

\begin{lemma}[R1: refinement computes the coarsest congruence]\label{lem:R1} \CLOSED\;
Let $M$ be a finite closed stub system. Then: (i) $M$ has a coarsest congruence
$\theta_{\max}$, and it equals the stable refinement partition $W(M)$ of
Definition~\ref{def:congruence}; (ii) $M$ is a core iff $W(M)$ is discrete; (iii) both
implementations of \texttt{simplify} (the solver's iterated sort-refinement and the
pruner's pair-elimination) compute $W(M)$, so \texttt{simplify} accepts $M$ iff $M$ is a
core.
\end{lemma}

\begin{proof}
(i) A partition is \emph{stable} iff whenever $x, y$ share a class, $c(x) = c(y)$ and
each of $\rho, \lambda, \mu, \gamma$ maps their classes into a common class. For
partitions into equivalence relations this stability is verbatim the compatibility
defining a congruence, because all four operations are (total, on a closed system)
functions: every congruence is a stable partition refining the $c$-fibers, and every
stable partition refining the $c$-fibers is a congruence. The coarsest stable partition
refining a given initial partition exists and is computed by iterated refinement
(standard partition-refinement theory; Cardon--Crochemore, Paige--Tarjan), which is
$W(M)$; hence $W(M) = \theta_{\max}$. (ii) is immediate from (i). (iii) The solver's
\texttt{simplify} is the textbook iterated refinement on the 6-tuple
$(c, \text{class}, \text{class}\circ\mu, \text{class}\circ\gamma,
\text{class}\circ\lambda, \text{class}\circ\rho)$ and converges to the coarsest stable
partition. The pruner's computes the greatest symmetric relation $R$ below the complete
relation that is compatible with colors and the four operations, by deleting violating
pairs to a fixpoint. Diagonal pairs never violate (their tests compare a pair to itself),
so $\Delta \subseteq R$; compatible symmetric relations containing $\Delta$ are closed
under composition of relations, and the transitive closure of $R$ is again compatible and
below the complete relation, so by maximality $R$ is transitive, hence an equivalence,
hence (by (i)) a congruence, so $R \subseteq \theta_{\max}$; conversely
$\theta_{\max}$ is a compatible symmetric relation below the complete one, so
$\theta_{\max} \subseteq R$ by maximality. Thus $R = W(M)$. (That the two code paths
implement these two computations faithfully is audit hook 6.)
\end{proof}

\begin{remark}[why Weisfeiler--Leman anxiety does not apply]
Color refinement is famously incomplete for isomorphism and automorphism detection on
general graphs. Stub systems are not general graphs: their relations are \emph{functions}
(deterministic transition structures), and on those, refinement computes the coarsest
congruence exactly; this is Myhill--Nerode/DFA-minimization logic, as
\texttt{algorithm.txt} itself hints by calling the check ``DFA minimization.'' Both
directions of the brief's obligation-4 worry (a rigid map wrongly rejected; two distinct
maps wrongly merged) are settled by R1 and P1 respectively.
\end{remark}

\begin{corollary}[R2: rejection drops nothing]\label{cor:R2} \CLOSED\;
If a closed valid $M$ fails the filter ($M$ not a core), the tiling $\dev(M)$ is still
enumerated: its canonical quotient $T/G(T)$, $T := \dev(M)$, is a core (B2b), is reached
independently (S1--S2), passes the filter (R1), and develops back to $T$ (B3(a)). Hence
discarding non-cores never removes a tiling from the catalog.
\end{corollary}

\begin{corollary}[R3: accepted solutions have exact uniformity]\label{cor:R3} \CLOSED\;
If $M$ is a core with $m$ vertices, then $G(\dev(M)) = H(M)$ and
$\Flags(\dev(M))/G(\dev(M)) \cong M$; in particular $\dev(M)$ is exactly $m$-uniform.
(Proof inside B3(c): a strictly larger symmetry group would induce a proper congruence
on $M$.)
\end{corollary}

\begin{lemma}[P1: pruner exactness on cores]\label{lem:P1} \CLOSED\;
Let $M, M'$ be finite connected closed stub systems that are cores. Then
\texttt{comparesolutions}$(M, M')$ returns true iff $M \cong M'$.
\end{lemma}

\begin{proof}
The procedure computes, by deleting violating pairs to a fixpoint, the greatest symmetric
relation $R$ on $X \sqcup X'$ that is contained in the initial relation
$\Delta \cup (X \times X') \cup (X' \times X)$ and compatible with colors and the four
operations. As in R1, $\Delta$ survives, and the procedure returns true iff some class of
$R$ is not a singleton, i.e.\ iff $R$ contains a cross pair.

($\Leftarrow$) If $\varphi : M \to M'$ is an isomorphism, the relation
$\Delta \cup \mathrm{graph}(\varphi) \cup \mathrm{graph}(\varphi^{-1})$ is symmetric,
compatible, and inside the initial relation; an easy induction shows the deletion process
never removes a pair of a compatible relation all of whose operation-image pairs it also
contains, so $R$ retains the cross pairs of $\varphi$ and the procedure returns true.

($\Rightarrow$) Suppose $R$ contains a cross pair and set
$R' := R \cap (X \times X')$, a nonempty relation compatible componentwise with colors
and operations (the fixpoint conditions transfer each cross pair's operation-images into
$R'$). The composite $S := R'^{\top} \circ R' \subseteq X' \times X'$ is symmetric,
compatible, and contains the diagonal of the range of $R'$; the equivalence it generates
together with $\Delta_{X'}$ is compatible (compatibility survives symmetric closure,
union with $\Delta$, and transitive closure, because the operations are functions), hence
is a congruence on $M'$; $M'$ a core forces it to be equality, so $R'$ is single-valued.
Symmetrically, $M$ a core forces $R'$ injective. The domain of $R'$ is closed under
$\rho, \lambda, \mu, \gamma$ (images of pairs are pairs), hence is all of $X$ by
connectedness; symmetrically $R'$ is surjective. A total, injective, surjective,
color- and operation-preserving relation that is single-valued is an isomorphism.

Both core hypotheses are supplied by the pipeline ordering invariant:
\texttt{simplify} runs before \texttt{comparesolutions} on every block, in file mode and
in \texttt{EU\_STREAM} mode alike. This ordering is a correctness precondition, not an
optimization; audit hook 6 pins it.
\end{proof}

\begin{lemma}[P2: bucketing never separates duplicates]\label{lem:P2} \CLOSED\;
The bucket key (letter-signature string, refinement fingerprint) is an isomorphism
invariant, so isomorphic solutions land in the same bucket and are compared; within a
bucket every new solution is compared against all kept ones. Hence, granting P1, exactly
one representative per isomorphism class survives.
\end{lemma}

\begin{proof}
The signature is the multiset of vertex letters; an isomorphism maps vertices to vertices
preserving $(\rho, \mu, c)$-structure, hence (by A6) preserves letters, hence the
multiset and its canonical string. The fingerprint is a hash of the sorted multiset of
per-stub colors after three rounds of refinement seeded by $c$; each round's recoloring
is defined from isomorphism-invariant data only, so isomorphic systems produce equal
multisets and equal hashes. (Hash collisions between non-isomorphic systems merely place
them in one bucket, where P1 separates them; nothing is lost.) The bucket therefore
contains every previously kept isomorph, and \texttt{compareToSeen} tests against all of
them.
\end{proof}

\begin{lemma}[P3: geometric cross-check via canonical form N]\label{lem:P3}
\MACHINE\ (optional, cheap)\;
For pruned $M \ne M'$ at the same $k$:
$N(\mathrm{symbol}(\dev(M))) \ne N(\mathrm{symbol}(\dev(M')))$, and for every
solver-emitted duplicate pair merged by the pruner the $N$-keys agree; $N$ is the proven
canonical form of \texttt{docs/canonical-form/}.
\end{lemma}

\begin{remark}[relation to canonical form N; answering the brief directly]
The pipeline computes no canonical form and does not need one: the pruner decides
isomorphism pairwise, and P1 proves that decision exact on cores. Canonical form $N$
lives at the geometric level and is related by the chain: combinatorial isomorphism of
cores $\iff$ congruence of developed tilings (B3) $\iff$ equal $N$-keys ($N$'s
Corollary 5.5). This makes $N$ an independent, already-proven arbiter that can
machine-audit every merge and non-merge on real runs (P3). If P1 had resisted closure,
replacing the pruner verdict by $N$-key equality at develop time was the designed escape
hatch; with P1 closed, P3 is belt and braces.
\end{remark}

%======================================================================
\section{Obligation 5: realizability (development)}\label{sec:realizability}

\begin{remark}[architecture: the direction bundle replaces orbifolds]\label{rem:bundle-arch}
Phase 1 planned this block through orbifold theory: realize $M$ as a compact flat
2-orbifold, quote developability (Thurston's orbifold chapter), pull back. That plan is
sound but imports machinery (orbifold gluings with mirror boundary, orbifold coverings)
whose fine print a referee must re-verify. The closed proofs below use a lighter route,
and it is the route \texttt{eu\_develop} literally computes: the \emph{direction bundle}
$B(M)$, the cover of $M$ that unfolds every site symmetry by remembering an edge
direction in $\Z/12$. Cone points and mirrors disappear upstairs: $B(M)$ assembles a
closed flat \emph{oriented} surface, necessarily a torus, and the only classical input
left is the classification of such surfaces. Orbifolds remain the right mental picture
of the quotient; they no longer occur in any proof.
\end{remark}

\begin{definition}[angle weights; property $(\angle)$]\label{def:angle}
For $n \in \mathcal{P}$ let $a(n) := 6 - 12/n \in \{2,3,4,5\}$, the interior angle of
the regular unit $n$-gon in units of $30^\circ$. A closed valid stub system has property
$(\angle)$ if every $\rho$-cycle $C$ satisfies
$s(C) := \sum_{x \in C} a(c(x)) \mid 12$. By A2(iv), every stub system over $\A$ has
$(\angle)$.
\end{definition}

\begin{lemma}[C0: congruence quotients preserve the hypotheses]\label{lem:C0} \CLOSED\;
If $M$ is finite, connected, closed, valid with $(\angle)$ and $\theta$ is a congruence
on $M$, then $M/\theta$ is finite, connected, closed, valid with $(\angle)$.
\end{lemma}

\begin{proof}
All structure descends (Definition~\ref{def:congruence}); finiteness, connectedness and
closedness are immediate. Validity: the quotient map $q$ intertwines $\omega$, so the
walk through $q(x)$ is the $q$-image of the walk through $x$; its minimal period $\ell'$
divides the original $\ell$, colors are constant as images of constants, and
$\ell' \mid \ell \mid n$ gives condition 2a; open chains do not occur. $(\angle)$: a
$\rho$-cycle $C'$ of $M/\theta$ is the image of a $\rho$-cycle $C$ with $|C'|$ dividing
$|C|$; the color word of $C$ is $|C'|$-periodic, so
$s(C') = s(C)\,|C'|/|C|$, which divides $s(C)$, hence $12$.
\end{proof}

\begin{definition}[the direction bundle]\label{def:bundle}
Let $M$ be finite, connected, closed, valid with $(\angle)$. Since $\mu$ normalizes
$\langle \rho, \gamma \rangle$, the $\langle \rho, \gamma \rangle$-orbits on $X$ are
either one ($X^+ := X$) or two swapped by $\mu$ ($X = X^+ \sqcup \mu X^+$; $M$ is then
called \emph{ultrachiral}, the solver's own term). Fix $X^+$ and a base stub
$x_0 \in X^+$. The \emph{direction bundle} is $B(M) := X^+ \times \Z/12$ with two
families of steps:
\[
\text{corner steps } (x,d) \to (\rho x,\ d + a(c(\rho x))), \qquad
\text{edge steps } (x,d) \to (\gamma x,\ d + 6).
\]
$B_0$ denotes the connected component of $(x_0, 0)$. Elements $(x,d)$ are \emph{darts};
\emph{bundle edges} are the pairs $\{(x,d), (\gamma x, d+6)\}$; the \emph{face walk} is
$\hat\omega :=$ (edge step, then corner step). The \emph{corner} between a dart and its
corner-step image is assigned the color $c(\rho x)$ and the angle $a(c(\rho x))$,
matching the corner convention of D1a.
\end{definition}

\begin{lemma}[C1: the bundle surface]\label{lem:C1} \CLOSED\;
(i) The corner-step cycle through $(x,d)$ projects onto the $\rho$-cycle $C \ni x$,
traverses it $12/s(C)$ times, and carries total angle weight exactly $12$.
(ii) Every face walk is a cycle; the walk through a dart of color $n$ has length exactly
$n$ and constant color $n$.
(iii) Every dart lies on exactly one face walk, and each bundle edge is traversed once
by each of the face walks through its two darts, in opposite directions.
(iv) Gluing one closed unit-edge regular $n$-gon per face walk of $B_0$ along the bundle
edges (isometrically on edges, reversing the boundary walk directions) yields a compact,
connected, flat, oriented surface $S(B_0)$.
\end{lemma}

\begin{proof}
(i) After $|C|$ corner steps the direction has advanced by $s(C)$; the cycle closes at
the least $t$ with $t\,s(C) \equiv 0 \pmod{12}$, namely $t = 12/s(C)$, an integer by
$(\angle)$; the total weight is $t \, s(C) = 12$.
(ii) One $\hat\omega$-step from a dart of color $n$ advances the direction by
$6 + a(n) \equiv -12/n \pmod{12}$ and projects to one $\omega$-step of $M$. The
underlying $M$-walk is a cycle of some length $\ell \mid n$ with constant color $n$
(validity), so the bundle walk can close only at multiples $t\ell$; the direction
returns first at $t\ell \cdot 12/n \equiv 0 \pmod{12}$, i.e.\ $t = n/\ell$: the bundle
walk is a cycle of length exactly $n$.
(iii) $\hat\omega$ is a bijection on darts, so its orbits partition them. The walk
through $(x,d)$ traverses the bundle edge $\{(x,d), (\gamma x, d{+}6)\}$ exactly when it
leaves $(x,d)$; the walk through the partner leaves $(\gamma x, d+6)$, crossing the same
edge with direction labels differing by $6$, i.e.\ oppositely.
(iv) The identification space of finitely many closed polygons under a complete
isometric pairing of boundary edges is a compact surface; its points are face-interior
points, edge-interior points (two flat half-disks), and vertex points. At a vertex
point, the corners glued around it are enumerated by alternately crossing a glued edge
and moving to the adjacent corner of the next polygon; under the conventions of
Definition~\ref{def:bundle} this alternation is the corner step (the corner between the
edges of $(x,d)$ and of its corner-step image is a corner of the polygon of color
$c(\rho x)$ glued along both), so by (i) the link is a single cycle of total angle $12$
units $= 2\pi$: a flat disk. Hence $S(B_0)$ is flat everywhere. Orienting every polygon
positively along its face walk, (iii) says each gluing reverses boundary orientations,
so the face orientations assemble to a global orientation. Compactness is finiteness of
$B_0$; connectedness is connectedness of $B_0$.
\end{proof}

\begin{lemma}[C2: classification]\label{lem:C2}
\CLASSICAL{Killing--Hopf; e.g.\ Wolf, \emph{Spaces of Constant Curvature}}\;
A closed flat oriented surface is isometric to $\C/\Lambda$ for a lattice
$\Lambda = \Z t_1 + \Z t_2 \subset \C$ of rank 2, acting by translations.
\end{lemma}

\begin{proof}[Proof note]
Classical. The elementary route: a compact flat surface is complete, so its universal
cover is the Euclidean plane (Killing--Hopf); deck transformations are fixed-point-free
isometries, orientation-preserving by orientability of the quotient; a
fixed-point-free orientation-preserving planar isometry is a translation; compactness
forces rank 2.
\end{proof}

\begin{definition}[development; the marked model]\label{def:dev}
Fix an isometry $S(B_0) \cong \C/\Lambda$ and let $p : \C \to S(B_0)$ be the covering
map. $\dev(M) := \hat T$ is the tiling of $\C$ whose tiles are the closures of the
connected lifts of the open faces of $S(B_0)$. Normalize the identification (rotating,
and conjugating if needed) so that the base dart $(x_0, 0)$ develops with edge direction
$0$ and corner steps rotate positively. A \emph{walk flag} of $\hat T$ is a flag reading
off a lifted face-walk traversal; $F^{\circ}(\hat T)$ denotes the walk flags (one of the
two orientation classes of $\Flags(\hat T)$). The \emph{marked model} is the bijection
\[
F^{\circ}(\hat T) \;\cong\; \{(b, z) : b \in B_0,\ z \in \mathrm{pos}(b) + \Lambda\},
\]
where $\mathrm{pos}$ is one fixed development of positions, under which corner steps fix
$z$ and the edge step at a dart with direction label $d$ adds $\zeta^d$, with
$\zeta := e^{i\pi/6}$. $\Phi : F^{\circ} \to X^+$ is the stub projection
$(b, z) \mapsto (\text{stub of } b)$.
\end{definition}

\begin{lemma}[C3a: the developed tiling]\label{lem:C3a} \CLOSED\;
$\hat T$ is an edge-to-edge, $\Lambda$-periodic tiling of $\C$ by unit-edge regular
polygons with colors in $\mathcal{P}$ and all edge directions on the $30^\circ$ grid:
the edge of every dart $(x,d)$ develops in direction $d \cdot 30^\circ$. The marked
model is well defined, and it is \emph{rigid}: a position assignment
$B_0 \to \C$ compatible with the step rules is determined by its value at one dart.
$\Phi$ commutes with $\rho$ and $\gamma$ and preserves colors. Tiles do not overlap and
leave no gaps by construction (the plane is the cover of an intrinsic gluing, not the
target of an immersion), which answers the brief's overlap worry and is why this block
will transfer to non-convex tiles unchanged. The exact area identity
$\sum \text{face areas} = |\det(t_1, t_2)|$ remains a runtime evidence gate, never a
proof step.
\end{lemma}

\begin{proof}
Faces of $S(B_0)$ have simply connected interiors, so each lifts isometrically to $\C$;
the closures are unit-edge regular $n$-gons; they have disjoint interiors and cover
$\C$ because the faces decompose $S(B_0)$ and $p$ is a covering; edge-to-edge because
the identifications pair full unit edges; $\Lambda$-periodic because $\Lambda$ is the
deck group. Directions: the flat metric on $\C/\Lambda$ has trivial holonomy, so every
edge of $S(B_0)$ has a well-defined direction in $\C$; assigning each dart its developed
direction, a corner step rotates by the interior angle $a \cdot 30^\circ$ positively
(normalization) and an edge step reverses ($+6$ units), matching the bundle labels; by
induction from the base dart along the connected $B_0$, the label $d$ \emph{is} the
direction in units. The same induction gives rigidity of the marked model: positions
propagate deterministically along corner steps (fixed) and edge steps (add $\zeta^d$),
and $B_0$ is connected. $\Phi$'s properties are read off the construction.
\end{proof}

\begin{lemma}[C3b: affine realization of bundle maps]\label{lem:C3b} \CLOSED\;
Let $M_1, M_2$ be as in Definition~\ref{def:bundle}, with developments
$\hat T_1, \hat T_2$ and marked models. Let $f : X_1^+ \to X_2^+$ be a surjection and
$\varepsilon \in \{+1, -1\}$, $h \in \Z/12$, such that
$\beta(x,d) := (f(x),\ \varepsilon d + h)$ maps steps of $B(M_1)$ to steps of $B(M_2)$
(for $\varepsilon = -1$: to reversed steps); explicitly, for $\varepsilon = +1$ this
holds iff $f$ commutes with $\rho, \gamma$ and preserves colors, and for
$\varepsilon = -1$ iff $f\rho = \rho^{-1}f$, $f\gamma = \gamma f$ and
$c \circ f = c \circ \rho$. Assume $\beta(B_0^{(1)}) \subseteq B_0^{(2)}$ (arrange by
composing with a direction shift). Then for every choice of a base pair there is a
unique isometry
\[
g(z) = \zeta^h z + w \quad (\varepsilon = +1), \qquad
g(z) = \zeta^h \bar z + w \quad (\varepsilon = -1)
\]
such that $(b, z) \mapsto (\beta b, g z)$ maps walk flags of $\hat T_1$ to walk flags of
$\hat T_2$; consequently $g(\hat T_1) = \hat T_2$.
\end{lemma}

\begin{proof}
Step compatibility of $(\beta, g)$: corner steps fix $z$ on both sides; the edge step at
$(x,d)$ adds $\zeta^d$, its $\beta$-image is the edge step at $(f x, \varepsilon d + h)$
adding $\zeta^{\varepsilon d + h}$, and indeed
$g(z + \zeta^d) = g(z) + \zeta^{h}\zeta^{d} = g(z) + \zeta^{d+h}$ for
$\varepsilon = 1$, respectively
$g(z + \zeta^d) = g(z) + \zeta^h \overline{\zeta^{d}} = g(z) + \zeta^{h-d}$ for
$\varepsilon = -1$. (The hypothesis on colors is what makes $\beta$ map corner steps
correctly: for $\varepsilon = -1$, the reversed corner step out of $f\rho x = \rho^{-1}fx$
adds $a(c(\rho\,\rho^{-1}fx)) = a(c(fx)) = a(c(\rho x))$, using
$c \circ f = c \circ \rho$; this is exactly the axiom $c\mu = c\rho$ when $f = \mu$.)
Pin $w$ by one base pair; by rigidity (C3a) the position assignment
$b \mapsto g(\mathrm{pos}_1(\beta^{-1}\text{-compatible base}))$ agrees with
$\mathrm{pos}_2$ up to the model's $\Lambda_2$-freedom, so $(\beta, g)$ maps every walk
flag of $\hat T_1$ to a walk flag of $\hat T_2$. Faces: by C1(ii) every face walk in
either bundle has length exactly its color, so $\beta$ restricted to a face walk is a
bijection onto a face walk of the same color, and $g$ maps the corresponding developed
tile isometrically onto a developed tile. Since $f$ is onto, every tile of $\hat T_2$
arises, so $g(\hat T_1) \subseteq \hat T_2$ tile-wise between two tilings, forcing
equality. Uniqueness: an isometry is determined by the image of one flag.
\end{proof}

\begin{lemma}[C3c: $M$ is the quotient of its development]\label{lem:C3c} \CLOSED\;
Define $\hat\Phi : \Flags(\hat T) \to X$ by $\hat\Phi := \Phi$ on walk flags and
$\hat\Phi := \mu \circ \Phi \circ \sigma_2$ on the other orientation class, and set
\[
H(M) := \{\, g \in \Isom(\C) \;:\; g(\hat T) = \hat T \text{ and }
\hat\Phi \circ g = \hat\Phi \,\}.
\]
Then $\hat\Phi$ is a surjective morphism of stub systems, $H(M) \le G(\hat T)$ acts
freely on flags with finitely many orbits, its orbits are exactly the $\hat\Phi$-fibers,
and $\hat\Phi$ descends to an isomorphism $\Flags(\hat T)/H(M) \cong M$. Moreover
$H(M)$ contains orientation-reversing elements iff $M$ is not ultrachiral.
\end{lemma}

\begin{proof}
\emph{$\hat\Phi$ is a morphism.} On walk flags this is C3a. On the other class, with
$\tilde y$ a walk flag: colors,
$c(\sigma_2\tilde y) = |{\rm face}(\tilde y)| = |{\rm face}(\sigma_1 \tilde y)| =
c_{\hat T}(\rho\tilde y) = c(\Phi\rho\tilde y) = c(\rho\Phi\tilde y) =
c(\mu\Phi\tilde y) = c(\hat\Phi \sigma_2\tilde y)$;
$\rho$: $\rho\sigma_2 = \sigma_2\rho^{-1}$, so
$\hat\Phi(\rho\,\sigma_2\tilde y) = \mu\Phi\rho^{-1}\tilde y = \rho\mu\Phi\tilde y
= \rho\,\hat\Phi(\sigma_2\tilde y)$; $\gamma$ commutes with $\sigma_2$; and
$\hat\Phi \circ \sigma_2 = \mu \circ \hat\Phi$ on both classes, so $\hat\Phi$
intertwines the mirrors. Surjectivity: onto $X^+$ by C3a; onto $\mu X^+$ (which is
$X^-$ or $X^+$) via the second clause.

\emph{Freeness.} An isometry fixing a flag is the identity (B0 applies to $\hat T$).

\emph{Orbits $=$ fibers.} $H(M)$ preserves fibers by definition. Transitivity, walk-flag
part: two walk flags over the same stub $x$ are $(b_1, z_1)$, $(b_2, z_2)$ with
$b_i = (x, d_i)$; the direction shift $\delta(y,e) := (y, e + d_2 - d_1)$ is a
$\beta$-map of C3b with $\varepsilon = 1$, $f = \mathrm{id}$, preserving $B_0$ (it maps
the component of $b_1$ to that of $b_2$, both $B_0$); realize it by
$g(z) = \zeta^{d_2-d_1} z + w$ pinned by $(b_1,z_1) \mapsto (b_2,z_2)$: then
$g \in H(M)$ and $g$ connects the two flags. Cross part: if $M$ is not ultrachiral,
$\mu$ maps $X^+$ to itself and satisfies the $\varepsilon = -1$ hypotheses of C3b (the
stub-system axioms verbatim), so after a direction-shift adjustment it is realized by
some orientation-reversing $g_-$ with $\hat\Phi \circ g_- = \hat\Phi$ (checked on walk
flags: $\hat\Phi(g_-\tilde y) = \mu\Phi(\sigma_2 g_- \tilde y) = \mu\,\mu\,\Phi\tilde y
= \Phi\tilde y$, since $\sigma_2 g_-$ is the flag realization of $\mu$); $g_-$ maps the
walk-flag part of a fiber to the other part, and $H(M) = H^+ \sqcup H^+ g_-$ acts
transitively on each fiber. If $M$ is ultrachiral, fibers over $X^+$ contain walk flags
only (the second clause of $\hat\Phi$ lands in $X^- = \mu X^+$), transitivity is the
walk-flag part, and no orientation-reversing $g$ can preserve $\hat\Phi$ (it would force
$\Phi$-values in $X^-$).

\emph{Quotient.} A surjective morphism whose fibers are the orbits of a freely acting
group descends to an isomorphism from the orbit space; finiteness of the orbit set is
$|X| < \infty$.
\end{proof}

\begin{lemma}[C3d: round trip and functoriality]\label{lem:C3d} \CLOSED\;
(i) If $T$ is a tiling and $H \le G(T)$ has finitely many flag orbits, then
$\dev(T/H)$ is congruent to $T$; normalizing $\dev(T/H) = T$, one has $H(T/H) = H$.
(ii) If $\theta$ is a congruence on $M$ and $q : M \to M/\theta$ the quotient, the
developments can be normalized so that $\dev(M) = \dev(M/\theta)$ as tilings of $\C$,
$\hat\Phi_{M/\theta} = q \circ \hat\Phi_M$, and $H(M) \le H(M/\theta)$.
(iii) If $M_1 \cong M_2$ then $\dev(M_1) \sim \dev(M_2)$.
\end{lemma}

\begin{proof}
(i) $T$'s own flags supply a step-compatible position assignment for $B_0(T/H)$: send a
walk flag of $T$ to (its image stub and direction, its vertex position); the step
equations hold in $T$ by the geometry of its tiles. Rotate $T$ so the base flag has
direction $0$; by rigidity (C3a) this assignment coincides with the marked model, so the
developed flag set of $\dev(T/H)$ equals the flag set of $T$: the tilings are equal.
$H(T/H) \supseteq H$ is immediate ($H$ preserves $T$ and the quotient labeling). For the
converse let $g \in H(T/H)$: $g$ preserves every fiber of
$\Flags(T) \to \Flags(T)/H$, i.e.\ every $H$-orbit; picking a flag $\tilde y$, there is
$h \in H$ with $g\tilde y = h\tilde y$, and $h^{-1}g$ fixes a flag, so $g = h \in H$ by
freeness (B0).
(ii) $\beta_q(x,d) := (q x, d)$ maps steps to steps (colors descend; C0 gives the
hypotheses on $M/\theta$), and bases can be chosen compatibly; C3b with
$\varepsilon = 1$, $h = 0$ realizes it by $g = \mathrm{id}$ once the two developments
are normalized with matching base positions: $\dev(M) = \dev(M/\theta)$ literally, and
$\hat\Phi_{M/\theta} = q \circ \hat\Phi_M$ by construction, whence every isometry
preserving $\hat\Phi_M$-fibers preserves the coarser $\hat\Phi_{M/\theta}$-fibers.
(iii) An isomorphism satisfies the $\varepsilon = 1$ hypotheses of C3b; the realized $g$
carries $\hat T_1$ to $\hat T_2$.
\end{proof}

\begin{remark}[the referee surface]\label{rem:referee}
Three places above compress convention verifications into a sentence; each is finite
bookkeeping, flagged here per the anti-fabrication discipline. (1) C1(iv): the
identification of the vertex-link alternation with the corner step (the corner between
the edges of a dart and of its corner-step image belongs to the polygon glued along
both). (2) C3a: the positivity convention for corner steps under the chosen orientation
(fixed by allowing a conjugation in the normalization). (3) C3c: the color computation
on the non-walk orientation class, written out in the proof but relying on the D1a
convention $c(y) = |{\rm face}(\sigma_2 y)|$. All three are checkable mechanically from
Definitions \ref{def:flags} and \ref{def:bundle}.
\end{remark}

\begin{lemma}[C4: \texttt{eu\_develop} computes the marked development]\label{lem:C4}
\CLOSED\ (abstract) + audit-facing\;
On input a finite connected closed valid $M$ with $(\angle)$, the BFS of
\texttt{eu\_develop} computes a position assignment for $B_0$ compatible with the step
rules (its \texttt{star} walk enumerates corner-step cycles; its glue step is the edge
step); it terminates within the guards; the recorded position discrepancies generate
exactly $\Lambda$; the integer elimination plus Lagrange--Gauss reduction output a basis
of $\Lambda$; and the seed reduction outputs exactly one representative position per
vertex of $\hat T$ modulo $\Lambda$. The developer therefore succeeds on every valid
input; in the pipeline it only ever runs on cores, for which $\Lambda$ is the full
translation lattice of $G(\dev(M))$ (R3), so the emitted cell is the true period cell.
\end{lemma}

\begin{proof}
The algorithm's keys are pairs (stub, direction) and its two moves are precisely the
corner and edge steps, so the explored key set is $B_0$ and the computed positions are a
step-compatible assignment; by rigidity (C3a) it is \emph{the} development up to the
base normalization. Termination and guards: a corner-step cycle has length equal to the
degree of the developed vertex, at most $6$, within the 60-step guard; the number of
keys is at most $12\,|X| \le 144\,\texttt{maxnum}$, within the 200k pop guard for every
reachable target (audit hook 8 converts both guards into assertions with these bounds).
Discrepancies: the BFS implicitly chooses a spanning tree of the step graph of $B_0$
(the tree of first placements); positions are propagated along tree steps, and every
non-tree step compares a propagated position against an existing one, recording the
difference, which lies in $\Lambda$ because both are positions of the same
$B_0$-point in the marked model. These differences generate $\Lambda$: the step graph of
$B_0$ is the vertex-edge graph of the torus tiling of $S(B_0)$ enriched with direction
data, its fundamental cycles surject onto $\pi_1$ of the surface (attaching the face
disks only adds relations), and the discrepancy of a cycle is exactly its
$\Lambda$-holonomy. The integer elimination (a Hermite-style column reduction) and
Lagrange--Gauss steps compute a basis of the generated lattice by standard integer
linear algebra, exactly (float enters only the reduction coefficient by rounding, never
position identity). The seed pass reduces the placed vertex positions modulo the basis
and dedupes to one representative per vertex class: with \emph{exact} coset reduction
(the canonical representative of Def.~2.2 of the $N$-note, two integer floor divisions)
this holds unconditionally. The shipped code instead reduces via float barycentric
coordinates, which agrees with the exact reduction whenever every developed coordinate
is bounded by $12\,|X| \le 144\,\texttt{maxnum}$ (each edge step adds a unit $\zeta^d$
and there are $\le 12\,|X|$ steps): the reachable range keeps every coordinate small (the
committed golden $k \le 6$ catalog has all reduced seed and lattice-basis coordinates
$\le 14$; instrumenting the pre-reduction developing map gives $\le 8$), against a float
break-point near $10^7$ where a translated congruent copy first mis-reduces --- five to
six orders of margin --- and \emph{fix-obligation FB-1} of the audit replaces the float
reduction by the exact one to remove the caveat entirely.
\end{proof}

%======================================================================
\section{Obligation 6: the planar--quotient bridge}\label{sec:bridge}

B0 is proved in Section~\ref{sec:model} (Lemma~\ref{lem:B0}).

\begin{lemma}[B1: $k$-uniform implies cocompact wallpaper, hence periodic]\label{lem:B1}
\CLOSED\ modulo citation pin (\CLASSICAL{Bieberbach; G\&S \S1.4/3.5 secondary})\;
If $T$ is $k$-uniform then $G(T)$ has at most $12k$ orbits on flags and acts cocompactly
on $\R^2$; being discrete (B0) and cocompact, it is one of the 17 wallpaper groups and
contains a rank-2 lattice of translations. Corollary: $k$-uniform implies doubly
periodic.
\end{lemma}

\begin{proof}
Flags at a vertex number $2d \le 12$ ($d \le 6$ since all angles are $\ge 60^\circ$), and
flags over one vertex orbit form at most $12$ flag orbits, so at most $12k$ in all.
Cocompactness: let $v_1, \dots, v_k$ represent the vertex orbits and let $R$ be the
circumradius of the dodecagon; every point of $\R^2$ lies in some tile, hence within $R$
of some vertex, hence within $R$ of $g v_i$ for some $g \in G(T)$, $i \le k$; so the
compact set $\bigcup_i \bar B(v_i, R)$ meets every orbit of points, i.e.\ the action is
cocompact. The classification of discrete cocompact planar isometry groups and the
existence of the translation lattice is Bieberbach's theorem in dimension 2 (citation to
pin; any standard crystallographic reference suffices).
\end{proof}

\begin{lemma}[B2a: the canonical quotient is a valid stub system over $\A$]\label{lem:B2a}
\CLOSED\;
For $[T] \in \T_k$, $M := T/G(T)$ is a finite connected closed valid stub system over
$\A$ with exactly $k$ vertices.
\end{lemma}

\begin{proof}
Finiteness, closedness, connectedness and the vertex correspondence are D1b with
$H = G(T)$ and B1 (finitely many flag orbits); the vertex count is $k$ by definition of
$k$-uniformity. Validity is L1(i). Over the alphabet: each vertex substructure is a fold
of its configuration word by a subgroup of the word's symmetry group (A2(i)--(ii)); the
configuration is one of A1's fourteen (O4 restricts sizes to $\mathcal{P}$, and the angle
equation holds at a tiling vertex); the alphabet lists one fold per conjugacy class
(A2(iii), certified against the table by A3), and conjugate folds are isomorphic, so the
substructure is isomorphic to an entry.
\end{proof}

\begin{lemma}[B2b: the canonical quotient is a core]\label{lem:B2b} \CLOSED\;
For a tiling $T$ with finitely many flag orbits, $T/G(T)$ is a core.
\end{lemma}

\begin{proof}
Let $\theta$ be a congruence on $M := T/G(T)$. By C3d(i) normalize
$\dev(M) = T$ with $H(M) = G(T)$. By C3d(ii), $H' := H(M/\theta)$ satisfies
$G(T) = H(M) \le H'$, and $H'$ preserves $\dev(M/\theta) = \dev(M) = T$, so
$H' \le G(T)$, forcing $H' = G(T)$. By C3c applied to $M/\theta$, the fibers of
$\hat\Phi_{M/\theta} = q \circ \hat\Phi_M$ are the $H'$-orbits $=$ the $G(T)$-orbits
$=$ the fibers of $\hat\Phi_M$; hence $q$ is injective and $\theta$ is equality.
\end{proof}

\begin{lemma}[B3: the Galois correspondence and the bijection]\label{lem:B3} \CLOSED\;
Let $T$ be a tiling and $H \le G(T)$ a subgroup with finitely many flag orbits;
normalize $\dev(T/H) = T$ (C3d(i)).
\begin{itemize}
\item[(a)] (Round trip.) $\Flags(\dev(M))/H(M) \cong M$ for every finite connected
  closed valid $M$ with $(\angle)$ (C3c), and $\dev(T/H) \sim T$ (C3d(i)).
\item[(b)] (Correspondence.) The assignments
  $H' \mapsto \theta_{H'} :=$ (the partition of $\Flags(T)/H$ induced by $H'$-orbits)
  and $\theta \mapsto H\!\left((T/H)/\theta\right)$ are mutually inverse,
  inclusion-preserving bijections between the intermediate groups
  $H \le H' \le G(T)$ and the congruences on $T/H$.
\item[(c)] (Bijection.) $M \mapsto [\dev(M)]$ and $[T'] \mapsto T'/G(T')$ are mutually
  inverse bijections between isomorphism classes of cores over $\A$ with $k$ vertices
  and $\T_k$. In particular every tiling has exactly one core, namely $T'/G(T')$.
\end{itemize}
\end{lemma}

\begin{proof}
(a) is C3c and C3d(i).

(b) \emph{Forward map well defined:} $H'$ normalizes nothing special but commutes with
the flag operations and preserves colors, so its orbit partition on $\Flags(T)$ descends
to a congruence on $\Flags(T)/H$ (classes are $H'$-orbits, unions of $H$-orbits since
$H \le H'$). \emph{Injective:} suppose $\theta_{H'} = \theta_{H''}$, i.e.\ $H'$ and
$H''$ have the same flag orbits. For $h' \in H'$ and a flag $\tilde y$ there is
$h'' \in H''$ with $h'\tilde y = h''\tilde y$; then $h''^{-1}h'$ fixes a flag, so
$h' = h'' \in H''$ (B0), giving $H' \subseteq H''$ and symmetrically equality.
\emph{Backward map lands in the interval:} for a congruence $\theta$ on $M := T/H$,
C3d(ii) gives $H \le H(M/\theta)$ and $H(M/\theta)$ preserves
$\dev(M/\theta) = \dev(M) = T$, so $H(M/\theta) \le G(T)$.
\emph{Backward then forward is the identity:} the fibers of
$\hat\Phi_{M/\theta} = q \circ \hat\Phi_M$ (C3d(ii)) are, by C3c applied to $M/\theta$,
the $H(M/\theta)$-orbits; their induced partition on $M = \Flags(T)/H$ is the fiber
partition of $q$, which is $\theta$.
\emph{Forward then backward is the identity:} for $H \le H' \le G(T)$, the quotient
$M/\theta_{H'}$ is canonically $\Flags(T)/H'$, and
$H(\Flags(T)/H') = H'$ by the argument of C3d(i) (an isometry preserving every
$H'$-orbit lies in $H'$, by freeness). \emph{Inclusion-preserving:} clear in both
directions.

(c) Well defined: if $M$ is a core with $k$ vertices, then applying (b) to
$T' := \dev(M)$, $H := H(M)$ (using $M \cong T'/H(M)$ from (a)): the congruence
$\theta_{G(T')}$ is proper iff $G(T') \supsetneq H(M)$; a core admits no proper
congruence, so $G(T') = H(M)$, the uniformity of $T'$ is the number of
$H(M)$-vertex-orbits $=$ the number of vertices of $M$, namely $k$; $T'$ is octagon-free
because its tiles carry colors in $\mathcal{P}$; so $[T'] \in \T_k$. Isomorphic cores
develop to congruent tilings (C3d(iii)), and congruent tilings have isomorphic canonical
quotients (D1b functoriality), so both maps are well defined on classes. Round trips:
$\dev(T'/G(T')) \sim T'$ is C3d(i); $T'/G(T') \cong M$ for $T' = \dev(M)$, $M$ a core,
is (a) together with $G(T') = H(M)$ just shown. B2a/B2b put $T'/G(T')$ in the right
codomain. Hence mutually inverse bijections, and the core of a tiling is unique.
\end{proof}

\subsection{Verdict on the brief's obligation-6 sharpening}\label{sec:verdict}

\textbf{Adopted, with corrections.} The brief's analysis is right on the two points that
matter, and both were verified rather than assumed:

\begin{enumerate}
\item \emph{No period bound is needed anywhere.} The enumerated object is the quotient
itself: a solution with $k$ vertex letters has at most $12k$ stubs, so finiteness of the
search (T1) is bare combinatorics, and no bound on the translation lattice of the
developed tiling enters T, S, C, or U. A period bound is what a \emph{geometric}
enumerator needs; demanding one here would import the mixed-type finiteness problem that
stalled the lattice method, for no benefit. Datta--Maity (arXiv:1705.05236) and
Kundu--Maity (arXiv:2111.15484) are demoted to context citations; independently, their
bounds are per-Archimedean-type and would not cover the mixed-type case anyway.
\item \emph{The split into 6a/6b is correct.} 6a (=B1) is classical: discreteness (B0)
plus cocompactness plus Bieberbach; ``$k$-uniform $\Rightarrow$ periodic'' falls out.
6b is the real content and is exactly B2 + B3 here.
\end{enumerate}

Corrections and refinements to the brief's formulation:
(i) discreteness of $G(T)$ must be established before Bieberbach applies (B0).
(ii) The quotient is in general a flat \emph{orbifold} with possible mirror boundary,
not a torus; but the closed proof does not run in the orbifold category at all: the
direction bundle (Remark~\ref{rem:bundle-arch}) passes to the translation-level cover,
where the object \emph{is} a torus, and the site symmetries return as the explicit
affine maps of C3b/C3c. The brief's instinct (``torus'') is thus right about the object
the proof should develop, wrong about the object the solver enumerates.
(iii) The flag bound is $12k$, not $\sim 24k$.
(iv) The brief's ``nondecreasing-vertex-type tie-break'' misdescribes the search: the
actual invariant is min-type rooting, which is weaker and easier to prove complete
(remark after S2).
(v) The brief's $k=1$ statement double-counts the octagon exclusion
(Remark~\ref{rem:hygiene}).
(vi) 6b should not be proved by importing Delaney--Dress realization theorems, though
the kinship deserves the remark in Section~\ref{sec:model}; self-contained 2D proofs are
shorter than the formal bridge to D-symbols would be.

%======================================================================
\section{Composition: proof of Theorem~\ref{thm:main} from the lemmas}\label{sec:compose}

\begin{proof}[Proof of Theorem~\ref{thm:main}, assuming the lemma blocks]
\textbf{(T)} is T1.

\textbf{(S)} Let $M \in P_k$. $M$ is a closed configuration built by the DFS, so it is
connected (every fresh vertex is glued to the existing component at birth), over $\A$ by
construction, and valid because the search vetoed every violation (S4, L1 monotonicity).
$M$ passed the rigidity filter, so $M$ is a core (R1). By the C-block (C1, C2, C3a),
$\dev(M)$ is an edge-to-edge tiling by unit regular polygons from $\mathcal{P}$, octagon
free by Convention~\ref{conv:scope}. By R3 its uniformity is exactly the vertex count of
$M$, which is $k$. Hence $[\dev(M)] \in \T_k$.

\textbf{(C)} Let $[T] \in \T_k$. By B1 and B2a/B2b, $M_0 := T/G(T)$ is a finite connected
closed valid core over $\A$ with exactly $k \le \texttt{maxnum}$ vertices. By S2 the run
roots a search at $M_0$'s minimal letter; by S1 that search visits a closed configuration
$M \cong M_0$; by R1 the filter accepts it; by S3 it is emitted; by P1--P2 the pruner
keeps exactly one representative of its isomorphism class, which is in $P_k$. By B3(a),
$\dev(M) \sim \dev(M_0) \sim T$.

\textbf{(U)} Let $M, M' \in P_k$ with $\dev(M) \sim \dev(M')$. Both are cores (R1), so by
B3(c) $M \cong T/G(T) \cong M'$ for $T := \dev(M)$. But the pruner never retains two
isomorphic solutions: they share a bucket (P2, using A6) and are compared and merged
(P1). Hence $M = M'$.

The bijection statement is (S) + (C) + (U).
\end{proof}

%======================================================================
\section{Lemma ledger}\label{sec:ledger}

\begingroup
\footnotesize
\begin{center}
\begin{tabular}{llll}
Lemma & Content & Tag & Depends on \\
\hline
D1a & flag calculus, dictionary, face walk & \textsc{closed} & --- \\
D1b & quotients; gluing-type dictionary & \textsc{closed} & D1a, B0 \\
B0 & flag rigidity, discreteness & \textsc{closed} & --- \\
O1 & octagon species: $\{4.8.8, 3.8.24\}$ & \textsc{closed} & --- \\
O2 & $3.8.24$ excluded & \textsc{closed} (+cite pin) & O1 \\
O3 & octagon forces $t1002$ & \textsc{closed} & O1, O2 \\
O4 & palette $\{3,4,6,12\}$, grid directions & \textsc{closed} (+cite pin) & --- \\
A1 & 14 configurations, 3 chiral & \textsc{closed} & --- \\
A2 & variants $=$ subgroup conjugacy classes & \textsc{closed} & D1a, B0 \\
A3 & 44 entries, three copies agree & \textsc{machine} (implemented) & A1, A2 \\
A4 & structural identities per entry & \textsc{machine} (implemented) & --- \\
A5 & ferk transversal / $\Aut$ certificate & \textsc{machine} (implemented) & A4 \\
A6 & entries pairwise non-isomorphic & \textsc{machine} (added; \textsc{pass}) & A3 \\
L1 & necessity + heredity of conditions & \textsc{closed} & D1a, D1b \\
L2 & mirror parity rule & \textsc{closed} & D1b \\
T1 & termination & \textsc{closed} & --- \\
S1 & guided descent (no-drop) & \textsc{closed} & L1, L2, A5, A6 \\
S2 & min-letter rooting; sharding partition & \textsc{closed} & S1 \\
S3 & emission gates & \textsc{closed} (abstract) + audit & --- \\
S4 & picked stub exists; tests = conditions & \textsc{closed} (abstract) + audit & --- \\
R1 & refinement $=$ coarsest congruence & \textsc{closed} & --- \\
R2 & rejection sound & \textsc{closed} (corollary) & R1, S1--S3, B2, B3 \\
R3 & exact uniformity of accepted cores & \textsc{closed} (corollary) & B3 \\
P1 & pruner exact on cores & \textsc{closed} & R1 \\
P2 & bucket keys are invariants & \textsc{closed} & A6 \\
P3 & N cross-check, both directions & \textsc{machine} (\textsc{pass}, $k \le 6$) & B3, N-note \\
C0 & quotients preserve valid/$(\angle)$ & \textsc{closed} & A2(iv) \\
C1 & bundle surface: flat, oriented, closed & \textsc{closed} & C0, A2(iv), L1 \\
C2 & closed flat oriented surface $= \C/\Lambda$ & \textsc{classical} (Killing--Hopf) & --- \\
C3a & developed tiling; grid; rigidity & \textsc{closed} & C1, C2 \\
C3b & affine realization of bundle maps & \textsc{closed} & C3a \\
C3c & $\Flags(\dev M)/H(M) \cong M$ & \textsc{closed} & C3b, B0 \\
C3d & round trip; functoriality & \textsc{closed} & C3b, C3c \\
C4 & \texttt{eu\_develop} $=$ marked development & \textsc{closed} + audit & C3a \\
B1 & cocompact wallpaper; periodicity & \textsc{closed} (Bieberbach) & B0 \\
B2a & $T/G(T)$ valid over $\A$, $k$ vtx & \textsc{closed} & D1b,A1,A2,A3,O4,L1,B1 \\
B2b & $T/G(T)$ is a core & \textsc{closed} & C3c, C3d \\
B3 & Galois correspondence; bijection & \textsc{closed} & C3c, C3d, B0 \\
\end{tabular}
\end{center}
\endgroup

After round 3: every proof row is \textsc{closed}; the \textsc{machine} rows are
discharged (A3--A5 audited and re-run, A6 added and passing, P3 passing in both
directions for $k \le 6$); the sole standing \textsc{classical} input is C2
(Killing--Hopf). The audit-facing halves of S3, S4, C4 are verified in deliverable B.
Deliverable A is complete as written, subject to the referee surface of
Remark~\ref{rem:referee}.

Round 3 (machine + audit), done: A6 certificate added to the generator; A3/A4/A5
audited against the lemma statements; P3 cross-check wired
(\texttt{scripts/p3-nkey-crosscheck.ts}); deliverable B written
(\texttt{audit-deliverable-B.md}); the $k{=}8$ exhibit root-caused and closed. An
adversarial-referee pass (2026-07-12) is folded in: it surfaced no defect in T/S/C/U for
the regular palette, and its confirmed minor findings (the O1 two-angle enumeration
above, the O3 Step-4 citation above, and the C4 magnitude caveat) are addressed in place.

%======================================================================
\section{Audit hooks (forward pointers to deliverable B)}\label{sec:audit}

The proof is about the abstract algorithm. Its contact points with the C++ are
enumerated here so deliverable B can check each; none is discharged yet.

\begin{enumerate}
\item \textbf{Tables.} A3/A4/A5 certificates: audit \texttt{gen\_alphabet.py --certify}
  against the lemma statements; confirm shipped binaries were built from certified
  tables; add the A6 certificate.
\item \textbf{Candidate enumeration} (S1): \texttt{extend} tries all free matching-parity
  stubs including self and mirror, and all letters $g \ge \mathrm{vertype}[0]$ with all
  \texttt{reps}; confirm the starred-stub normalization of \texttt{firstfree} loses no
  candidate.
\item \textbf{Rooting and sharding} (S2): \texttt{initex} covers every letter;
  \texttt{EU\_SHARD\_N/W} partitions roots; union-equals-sequential.
\item \textbf{Emission} (S3): \texttt{seen} $= 0$; noncounting budget inert on the
  regular palette; \texttt{EU\_STREAM} changes the sink only; no other cap in the write
  path.
\item \textbf{Local tests} (S4): \texttt{checkpart} against conditions 1/2a/2b;
  \texttt{writecycle}/\texttt{mincycle} free-stub selection.
\item \textbf{Pipeline ordering invariant} (R1(iii), P1): \texttt{simplify} precedes
  \texttt{comparesolutions} on every path; the two refinement implementations compute
  the stated fixpoints.
\item \textbf{Bucketing} (P2): key is (signature line, fingerprint); buckets never cross
  $k$; \texttt{EU\_KONLY} semantics.
\item \textbf{Developer} (C4): star-walk and BFS guards replaced by (or asserted
  against) proven bounds; float confined to lattice reduction and coset rounding; the
  \texttt{--sign} convention; area certificate logged as evidence, not proof.
\item \textbf{The $k=8/9$ exhibit}: \textsc{closed} in deliverable B
  (\texttt{audit-deliverable-B.md}, \S4, same directory). Root cause: the
  Python reference solver's hand-computed attachment bound for the letter
  $(4,4,4,4)A2$ tried one dart instead of two, missing the starred Aut-orbit; exactly
  one $k=8$ tiling (species partition $3{+}3{+}1{+}1$, group $p4m$) is reachable only
  through a gluing at that dart and was silently dropped. A one-entry fix restores
  $2850$. This is Lemma A5's failure class realized in the wild, and the A5/A6
  certificates are exactly the checks that exclude it from the audited engine.
\end{enumerate}

%======================================================================
\begin{thebibliography}{9}
\bibitem{GS} B.~Gr\"unbaum, G.~C.~Shephard, \emph{Tilings and Patterns}, W.~H.~Freeman,
1987. Pinned against the copy in \texttt{resources/papers}: equation (2.1.2) (the angle
equation); Table 2.1.1, pp.~59--61 (17 species, 21 types; the six double-starred species
$3.7.42$, $3.8.24$, $3.9.18$, $3.10.15$, $4.5.20$, $5.5.10$ admit no vertex in any
edge-to-edge tiling, p.~59); Statement 2.1.3 (the 11 uniform tilings); enantiomorphy of
$(3^4.6)$, p.~59.
\bibitem{Bieberbach} L.~Bieberbach, \"Uber die Bewegungsgruppen der Euklidischen
R\"aume I, \emph{Math.\ Ann.} 70 (1911) 297--336. (Cocompact discrete planar groups; any
standard crystallographic-groups reference may substitute.)
\bibitem{Wolf} J.~A.~Wolf, \emph{Spaces of Constant Curvature}, 6th ed., AMS Chelsea,
2011. (Space-form classification; used for C2. Any source for Killing--Hopf plus the
fixed-point-free classification of planar isometries substitutes.)
\bibitem{Thurston} W.~P.~Thurston, \emph{The Geometry and Topology of Three-Manifolds},
Princeton lecture notes, ch.~13 (orbifolds). (Context only since round 2: the orbifold
route was replaced by the direction bundle, Remark in Section 8.)
\bibitem{Dress} A.~W.~M.~Dress, Presentations of discrete groups, acting on simply
connected manifolds, \emph{Adv.\ Math.} 63 (1987) 196--212; A.~W.~M.~Dress, D.~Huson,
On tilings of the plane, \emph{Geometriae Dedicata} 24 (1987) 295--310. (Context: the
D-symbol correspondence; not imported.)
\bibitem{CC} A.~Cardon, M.~Crochemore, Partitioning a graph in $O(|A|\log_2|V|)$,
\emph{Theoret.\ Comput.\ Sci.} 19 (1982) 85--98; R.~Paige, R.~E.~Tarjan, Three partition
refinement algorithms, \emph{SIAM J.\ Comput.} 16 (1987) 973--989. (Partition
refinement computes the coarsest stable partition; used in R1.)
\bibitem{Ctrnact} M.~\v{C}trn\'act, tiling solver/pruner algorithm description and
implementation, \texttt{tools/ctrnact-oracle/reference/} (\texttt{algorithm.txt},
\texttt{eu\_solver.orig.cpp}), 2021--2022.
\bibitem{N} TilingAtlas working note, \emph{A canonical form for integer symbols of
periodic tilings by regular polygons}, \texttt{docs/canonical-form/}, July 2026.
\bibitem{DM} B.~Datta, D.~Maity, arXiv:1705.05236; A.~K.~Kundu, D.~Maity,
arXiv:2111.15484. (Context only; demoted per Section~\ref{sec:verdict}.)
\end{thebibliography}

\end{document}
